% ============================================================
%  يُترجَم بـ XeLaTeX  (في Overleaf: Menu → Compiler → XeLaTeX)
%  الخط: Amiri (متوفر في Overleaf، أو حمّله من github.com/aliftype/amiri)
% ============================================================
\documentclass[12pt,a4paper]{article}
\usepackage[a4paper,margin=1.7cm,top=2.3cm,bottom=2.9cm,headheight=15pt,footskip=1.5cm]{geometry}
\usepackage{amsmath,amssymb}
\usepackage{icomma}
\usepackage[english,bidi=default,layout=lists]{babel}
\babelprovide[import,main]{arabic}
\babelfont{rm}[Scale=1.08]{Amiri}
\babelfont{sf}[Scale=1.08]{Amiri}
\usepackage[most]{tcolorbox}
\usepackage{tikz}
\usetikzlibrary{arrows.meta,calc,shadings}
\usepackage{array}
\usepackage{colortbl}
\usepackage{fancyhdr}
\usepackage{titlesec}
\usepackage{enumitem}
\usepackage{collcell}
\usepackage{needspace}
\usepackage{etoolbox}
\newcolumntype{C}{>{\collectcell\ar}c<{\endcollectcell}}
\pretocmd{\section}{\Needspace{15\baselineskip}}{}{}
\pretocmd{\subsection}{\Needspace{11\baselineskip}}{}{}

\definecolor{cDef}{HTML}{1F4E9C}
\definecolor{cThm}{HTML}{B3261E}
\definecolor{cRem}{HTML}{2E7D32}
\definecolor{cEx}{HTML}{C76B00}
\definecolor{cMeth}{HTML}{6A1B9A}
\definecolor{cSol}{HTML}{00695C}
\definecolor{cSec}{HTML}{1F3A68}

\titleformat{\section}{\Large\bfseries\color{cSec}}{\thesection.}{0.6em}{}[{\color{cSec!60}\titlerule[1pt]}]
\titleformat{\subsection}{\large\bfseries\color{cThm!85!black}}{\thesubsection.}{0.5em}{}
\titlespacing*{\section}{0pt}{14pt}{8pt}
\titlespacing*{\subsection}{0pt}{10pt}{5pt}
\setlength{\parindent}{0pt}
\setlength{\parskip}{4pt}
\setlist{itemsep=2pt,topsep=3pt}

\newtcolorbox{arbox}[3][]{enhanced,
  colback=#2!5,colframe=#2!80!black,arc=2mm,boxrule=0.8pt,
  left=6pt,right=6pt,top=8pt,bottom=5pt,
  fonttitle=\bfseries,coltitle=white,
  title={\foreignlanguage{arabic}{#3}},
  attach boxed title to top right={xshift=-5mm,yshift=-3mm},
  boxed title style={colback=#2!80!black,arc=1.5mm,boxrule=0pt},
  before upper={\selectlanguage{arabic}},
  before skip=8pt,after skip=8pt,breakable,#1}
\newenvironment{BOX}[2]{\begin{otherlanguage}{english}\begin{arbox}{#1}{#2}}{\end{arbox}\end{otherlanguage}}
\newenvironment{tarif}[1][تعريف]{\begin{BOX}{cDef}{#1}}{\end{BOX}}
\newenvironment{khassia}[1][خاصية]{\begin{BOX}{cThm}{#1}}{\end{BOX}}
\newenvironment{molahada}[1][ملاحظة]{\begin{BOX}{cRem}{#1}}{\end{BOX}}
\newenvironment{mithal}[1][مثال]{\begin{BOX}{cEx}{#1}}{\end{BOX}}
\newenvironment{tariqa}[1][طريقة]{\begin{BOX}{cMeth}{#1}}{\end{BOX}}

\newcounter{exo}
\newenvironment{tamrin}[1]{\stepcounter{exo}\begin{BOX}{cSec}{التمرين \theexo: #1}}{\end{BOX}}
\newenvironment{hal}{\Needspace{12\baselineskip}\begin{otherlanguage}{english}\begin{arbox}[breakable]{cSol}{حل التمرين \theexo}}{\end{arbox}\end{otherlanguage}}

\newenvironment{LTR}{\begin{otherlanguage}{english}}{\end{otherlanguage}}
\newcommand{\ar}[1]{\foreignlanguage{arabic}{#1}}

\renewcommand{\dots}{\mathinner{\ldotp\ldotp\ldotp}}
\renewcommand{\ldots}{\mathinner{\ldotp\ldotp\ldotp}}
\renewcommand{\cdots}{\mathinner{\cdotp\cdotp\cdotp}}
\newcommand{\N}{\mathbb{N}}
\newcommand{\Z}{\mathbb{Z}}
\newcommand{\D}{\mathbb{D}}
\newcommand{\Q}{\mathbb{Q}}
\newcommand{\R}{\mathbb{R}}
\newcommand{\pgcd}{\operatorname{PGCD}}
\newcommand{\ppcm}{\operatorname{PPCM}}

% \NLseg{a}{b}{[}{]} ; \NLright{a}{[} ; \NLleft{b}{]}  (interval on a number line)
\newcommand{\NLseg}[4]{\begin{LTR}\begin{tikzpicture}[baseline=-0.6ex]
 \draw[-{Stealth[length=2mm]},gray!70](0,0)--(3.4,0);
 \draw[line width=2pt,cThm!80!black](0.7,0)--(2.7,0);
 \node[cThm!80!black,font=\LARGE] at (0.7,0){$#3$};
 \node[cThm!80!black,font=\LARGE] at (2.7,0){$#4$};
 \node[below=3pt,font=\footnotesize] at (0.7,0){$#1$};
 \node[below=3pt,font=\footnotesize] at (2.7,0){$#2$};
\end{tikzpicture}\end{LTR}}
\newcommand{\NLright}[2]{\begin{LTR}\begin{tikzpicture}[baseline=-0.6ex]
 \draw[-{Stealth[length=2mm]},gray!70](0,0)--(3.4,0);
 \draw[line width=2pt,cThm!80!black](1.3,0)--(3.3,0);
 \node[cThm!80!black,font=\LARGE] at (1.3,0){$#2$};
 \node[below=3pt,font=\footnotesize] at (1.3,0){$#1$};
\end{tikzpicture}\end{LTR}}
\newcommand{\NLleft}[2]{\begin{LTR}\begin{tikzpicture}[baseline=-0.6ex]
 \draw[-{Stealth[length=2mm]},gray!70](0,0)--(3.4,0);
 \draw[line width=2pt,cThm!80!black](0.05,0)--(2.1,0);
 \node[cThm!80!black,font=\LARGE] at (2.1,0){$#2$};
 \node[below=3pt,font=\footnotesize] at (2.1,0){$#1$};
\end{tikzpicture}\end{LTR}}

% ---------- social-media icons (pure TikZ) ----------
\definecolor{cFB}{HTML}{1877F2}
\definecolor{cTG}{HTML}{229ED9}
\definecolor{cIG1}{HTML}{FEDA75}
\definecolor{cIG2}{HTML}{D62976}
\definecolor{cIG3}{HTML}{4F5BD5}
\newcommand{\fbicon}{\begin{LTR}\begin{tikzpicture}[baseline=-0.75ex,x=1em,y=1em]
 \fill[cFB] (0,0) circle (0.55);
 \node[white,inner sep=0,font=\usefont{TU}{lmss}{bx}{n}\small] at (0.04,-0.02) {f};
\end{tikzpicture}\end{LTR}}
\newcommand{\igicon}{\begin{LTR}\begin{tikzpicture}[baseline=-0.75ex,x=1em,y=1em]
 \shade[lower left=cIG1,lower right=cIG2,upper left=cIG2,upper right=cIG3,rounded corners=0.28em] (-0.55,-0.55) rectangle (0.55,0.55);
 \draw[white,line width=0.09em,rounded corners=0.18em] (-0.36,-0.36) rectangle (0.36,0.36);
 \draw[white,line width=0.09em] (0,0) circle (0.17);
 \fill[white] (0.22,0.22) circle (0.05);
\end{tikzpicture}\end{LTR}}
\newcommand{\tgicon}{\begin{LTR}\begin{tikzpicture}[baseline=-0.75ex,x=1em,y=1em]
 \fill[cTG] (0,0) circle (0.55);
 \fill[white] (-0.33,0.0) -- (0.32,0.27) -- (0.2,-0.3) -- (0.0,-0.15) -- (-0.08,-0.28) -- (-0.08,-0.08) -- cycle;
 \fill[cTG!70!black] (-0.08,-0.08) -- (0.25,0.2) -- (0.0,-0.15) -- (-0.08,-0.28) -- cycle;
\end{tikzpicture}\end{LTR}}
% page number: black box, white two-digit number, at the right of the social row
\newcommand{\PageBox}{\begin{LTR}\tikz[baseline=(p.base)]\node[fill=black,text=white,inner sep=2.5pt,minimum width=1.6em,minimum height=1.3em,font=\usefont{TU}{lmss}{bx}{n}\small](p){\ifnum\value{page}<10 0\fi\arabic{page}};\end{LTR}}
% visual left-to-right order: Telegram | Instagram | Facebook | page box
\newcommand{\SocialRow}{\begin{otherlanguage}{english}%
 \ar{\textbf{الأستاذ سوايسية هشام للرياضيات}}\ \tgicon\hfill
 \textbf{prof\_hicham\_math}\ \igicon\hfill
 \ar{\textbf{الأستاذ سوايسية هشام للرياضيات}}\ \fbicon\hspace{2.2em}\PageBox\end{otherlanguage}}

\pagestyle{fancy}
\fancyhf{}
\renewcommand{\headrulewidth}{0.6pt}
\renewcommand{\footrulewidth}{0.4pt}
% \SetHeader{right head: document title}{left head: level -- stream}
\newcommand{\SetHeader}[2]{%
 \fancyhead[C]{\parbox{\textwidth}{\ar{#1\hfill #2}}}%
 \fancyfoot[C]{\parbox{\textwidth}{\small\SocialRow}}}
\usepackage{tkz-tab}
\usepackage{multicol}
% variation / sign table (LTR)
\newenvironment{TAB}{\par\begin{center}\begin{LTR}\begin{tikzpicture}[scale=0.95,every node/.style={font=\small}]}{\end{tikzpicture}\end{LTR}\end{center}}
% grid + axes, clipped: \begin{GR}{xmin}{xmax}{ymin}{ymax}{unit cm}
\newenvironment{GR}[5]{\par\begin{center}\begin{LTR}\begin{tikzpicture}[x=#5cm,y=#5cm,font=\footnotesize]
 \draw[gray!25,very thin] (#1,#3) grid (#2,#4);
 \draw[-{Stealth[length=2mm]}] (#1,0)--(#2+0.3,0) node[below left]{$x$};
 \draw[-{Stealth[length=2mm]}] (0,#3)--(0,#4+0.3) node[below left]{$y$};
 \node[below left] at (0,0){$O$};
 \begin{scope}\clip (#1,#3) rectangle (#2,#4);}
{\end{scope}\end{tikzpicture}\end{LTR}\end{center}}
\newcommand{\Ticks}[2]{\foreach \t in {#1} \node[below,font=\scriptsize] at (\t,0){$\t$}; \foreach \t in {#2} \node[left,font=\scriptsize] at (0,\t){$\t$};}
\newcommand{\bx}[1]{\ensuremath{\colorbox{cDef!12}{$\displaystyle #1$}}}
\newcommand{\TwoCol}[4]{\par\begin{LTR}\noindent\begin{minipage}[t]{#1\textwidth}\selectlanguage{arabic}#2\end{minipage}\hfill\begin{minipage}[t]{#3\textwidth}\selectlanguage{arabic}#4\end{minipage}\end{LTR}\par}
\newcommand{\RL}[4]{\TwoCol{#3}{#4}{#1}{#2}}  % \RL{w1}{right content}{w2}{left content}
\newcommand{\ok}{$\textcolor{cRem}{\checkmark}$}
\newcommand{\limn}{\displaystyle\lim_{n\to+\infty}}
% sub-questions, each on its own line
\newenvironment{sq}{\begin{itemize}[leftmargin=2.4em,itemsep=3pt]}{\end{itemize}}
\newcommand{\qa}{\item[(أ)]}\newcommand{\qb}{\item[(ب)]}\newcommand{\qc}{\item[(جـ)]}\newcommand{\qd}{\item[(د)]}\newcommand{\qe}{\item[(هـ)]}
\newcommand{\pa}{\par\smallskip\textbf{(أ)} }\newcommand{\pb}{\par\smallskip\textbf{(ب)} }\newcommand{\pc}{\par\smallskip\textbf{(جـ)} }\newcommand{\pd}{\par\smallskip\textbf{(د)} }\newcommand{\pe}{\par\smallskip\textbf{(هـ)} }
\newcommand{\QQ}[1]{\par\smallskip\textbf{#1.} }
\newenvironment{tmr}{\stepcounter{exo}\begin{BOX}{cSec}{تمرين \theexo}}{\end{BOX}}
\newcommand{\ABC}[3]{\\[2pt]\begin{tabular}{p{0.28\textwidth}p{0.28\textwidth}p{0.28\textwidth}}(جـ) \ #3 & (ب) \ #2 & (أ) \ #1\end{tabular}}
\pretocmd{\section}{\Needspace{24\baselineskip}}{}{}
\newtcbox{\bxtbox}{on line,colback=cDef!12,colframe=cDef!12,boxrule=0pt,arc=1pt,boxsep=0pt,left=2pt,right=2pt,top=1pt,bottom=1pt}
\newcommand{\bxt}[1]{\foreignlanguage{english}{\bxtbox{\ar{#1}}}}
\SetHeader{\small\bfseries\color{cSec}تمارين البكالوريا: الدالة الأسية}{\small\color{cSec}السنة الثالثة ثانوي -- الشعب العلمية}
\newenvironment{FIG}{\par\begin{center}\begin{LTR}}{\end{LTR}\end{center}}
\begin{document}
\begin{LTR}
\begin{tcolorbox}[enhanced,colback=cSec,colframe=cSec,arc=3mm,halign=center,boxrule=0pt,top=12pt,bottom=12pt,before upper={\selectlanguage{arabic}}]
{\color{white}\fontsize{26}{32}\selectfont\bfseries تمارين البكالوريا في الدالة الأسية}\\[8pt]
{\color{white}\large -- بدون حلول (للتدرّب) -- حتى دورة 2026}\\[5pt]
{\color{white}\large علوم تجريبية -- تقني رياضي -- رياضيات}
\end{tcolorbox}
\end{LTR}
\section{تمارين البكالوريا -- شعبة علوم تجريبية}
% ===================== بكالوريا 2026 -- الموضوع الثاني =====================
\begin{tamrin}{بكالوريا 2026 -- الموضوع الثاني -- علوم تجريبية}
$f$ الدالة المعرفة على $\R^{*}$ بـ: $f(x)=1+\dfrac{5-4e^{x}}{e^{2x}-1}$، و $(C_f)$ تمثيلها البياني في المستوي المنسوب إلى المعلم المتعامد والمتجانس $(O;\vec i,\vec j)$.
\QQ{1} أحسب النهايات: $\lim\limits_{x\to-\infty}f(x)$، $\lim\limits_{x\to0^{+}}f(x)$، $\lim\limits_{x\to0^{-}}f(x)$ و $\lim\limits_{x\to+\infty}f(x)$، ثم فسّر النتائج هندسيا.
\QQ{2}
\begin{sq}
\qa بيّن أنه من أجل كل $x$ من $\R^{*}$: $f'(x)=\dfrac{2e^{x}\left(e^{x}-2\right)\left(2e^{x}-1\right)}{\left(e^{2x}-1\right)^{2}}$.
\qb استنتج اتجاه تغيّر الدالة $f$ ثم شكّل جدول تغيّراتها.
\end{sq}
\QQ{3} بيّن أنه من أجل كل $x$ من $\R^{*}$: $f(-x)+f(x)=-3$، وفسّر النتيجة هندسيا، ثم أرسم المنحنى $(C_f)$.
\QQ{4} تحقّق أنه من أجل كل $x$ من $\R^{*}$: $f(x)=-4+\dfrac{1}{2}\left(\dfrac{9e^{x}}{1+e^{x}}-\dfrac{e^{x}}{1-e^{x}}\right)$.
\QQ{5} $g$ الدالة المعرفة على $]1;+\infty[$ بـ: $g(x)=f(\ln x)$.
تحقّق أنه من أجل كل $x>1$: $g(x)=\dfrac{(x-2)^{2}}{x^{2}-1}$، وحدّد اتجاه تغيّر الدالة $g$ ثم شكّل جدول تغيّراتها.
\end{tamrin}

\begin{molahada}[تصحيح في نص التمرين]
في السؤال 4 كُتب المقام $1+e^{2}$ وهو خطأ مطبعي؛ الصواب $1+e^{x}$ (بالمقام $1+e^{2}$ لا تتحقق المساواة). صُحّح في النص أعلاه.
\end{molahada}

% ===================== بكالوريا 2025 -- الموضوع الأول =====================
\begin{tamrin}{بكالوريا 2025 -- الموضوع الأول -- علوم تجريبية}
$f$ الدالة المعرفة على $\R$ بـ: $f(x)=e^{2x}-e^{x}-x-2$، و $(C_f)$ تمثيلها البياني في المستوي المنسوب إلى المعلم المتعامد والمتجانس $(O;\vec i,\vec j)$ (وحدة الطول $2cm$).
\QQ{1}
\begin{sq}
\qa أحسب $\lim\limits_{x\to-\infty}f(x)$ ثم بيّن أن $\lim\limits_{x\to+\infty}f(x)=+\infty$.
\qb بيّن أن المستقيم $(\Delta)$ ذا المعادلة $y=-x-2$ مقارب مائل لـ $(C_f)$ عند $-\infty$.
\qc أدرس الوضع النسبي للمنحنى $(C_f)$ والمستقيم $(\Delta)$.
\end{sq}
\QQ{2}
\begin{sq}
\qa بيّن أنه من أجل كل عدد حقيقي $x$: $f'(x)=\left(e^{x}-1\right)\left(2e^{x}+1\right)$.
\qb استنتج اتجاه تغيّر الدالة $f$ ثم شكّل جدول تغيّراتها.
\end{sq}
\QQ{3} بيّن أن المنحنى $(C_f)$ يقبل مماسا $(T)$ موازيا للمستقيم $(\Delta)$، يُطلب تعيين معادلة له.
\QQ{4} بيّن أن المعادلة $f(x)=0$ تقبل حلّين فقط $\alpha$ و $\beta$، ثم تحقّق أن: $-2{,}2<\alpha<-2{,}1$ و $0{,}8<\beta<0{,}9$.
\QQ{5} بيّن أن المنحنى $(C_f)$ يقبل نقطة انعطاف يُطلب تعيين إحداثيّيها.
\QQ{6}
\begin{sq}
\qa أرسم كلا من $(\Delta)$، $(T)$ و $(C_f)$.
\qb ناقش بيانيا، وحسب قيم الوسيط الحقيقي $m$، عدد حلول المعادلة: $e^{2x}-e^{x}-m-2=0$.
\end{sq}
\end{tamrin}

% ===================== بكالوريا 2024 -- الموضوع الأول =====================
\begin{tamrin}{بكالوريا 2024 -- الموضوع الأول -- علوم تجريبية}
\textbf{I)} يمثل الجدول التالي تغيّرات الدالة $g$ المعرفة على $\R$ بـ: $g(x)=xe^{-x+1}-2$.
\begin{TAB}
\tkzTabInit[lgt=2,espcl=3]{$x$/1,$g'(x)$/1,$g(x)$/2}{$-\infty$,$1$,$+\infty$}
\tkzTabLine{,+,z,-,}
\tkzTabVar{-/$-\infty$,+/$g(1)$,-/$-2$}
\end{TAB}
أحسب $g(1)$، ثم استنتج إشارة $g(x)$.

\textbf{II)} $f$ الدالة المعرفة على $\R$ بـ: $f(x)=-2x+3-xe^{-x+1}$، و $(C_f)$ تمثيلها البياني في المستوي المنسوب إلى المعلم المتعامد والمتجانس $(O;\vec i,\vec j)$ (وحدة الطول $2cm$).
\QQ{1}
\begin{sq}
\qa أحسب $\lim\limits_{x\to-\infty}f(x)$ و $\lim\limits_{x\to+\infty}f(x)$.
\qb بيّن أن المستقيم $(\Delta)$ ذا المعادلة $y=-2x+3$ مقارب مائل للمنحنى $(C_f)$ عند $+\infty$، ثم أدرس الوضع النسبي للمنحنى $(C_f)$ والمستقيم $(\Delta)$.
\end{sq}
\QQ{2}
\begin{sq}
\qa بيّن أنه من أجل كل عدد حقيقي $x$: $f'(x)=g(x)-e^{-x+1}$.
\qb استنتج اتجاه تغيّر الدالة $f$، ثم شكّل جدول تغيّراتها.
\end{sq}
\QQ{3} بيّن أن $(C_f)$ يقبل مماسا $(T)$ موازيا لـ $(\Delta)$، يُطلب تعيين معادلة له.
\QQ{4}
\begin{sq}
\qa أرسم $(T)$، $(\Delta)$ و $(C_f)$.
\qb عيّن بيانيا قيم الوسيط الحقيقي $m$ التي من أجلها تقبل المعادلة $f(x)=-2x+m$ حلّين مختلفين.
\end{sq}
\end{tamrin}

% ===================== بكالوريا 2023 -- الموضوع الأول =====================
\begin{tamrin}{بكالوريا 2023 -- الموضوع الأول -- علوم تجريبية}
\textbf{I)} $(\Gamma)$ التمثيل البياني للدالة المعرفة على $\R$ بـ: $x\mapsto(2x-1)e^{2x}$، و $(D)$ المستقيم ذو المعادلة $y=1$، و $\alpha$ فاصلة نقطة تقاطع $(\Gamma)$ و $(D)$ (لاحظ الشكل).
\begin{GR}{-3.5}{1.6}{-1.5}{3.5}{0.9}
\Ticks{-3,-2,-1,1}{-1,1,2,3}
\draw[cThm,thick](-3.5,1)--(1.6,1) node[pos=0.1,above]{$(D)$};
\draw[cDef,very thick,domain=-3.5:1.2,samples=150] plot(\x,{(2*\x-1)*exp(2*\x)});
\fill(0.639,1) circle(1.5pt); \draw[densely dotted](0.639,1)--(0.639,0) node[below right]{$\alpha$};
\node[cDef] at (-1.6,-0.6){$(\Gamma)$};
\end{GR}
\QQ{1} بقراءة بيانية، حدّد وضعية $(\Gamma)$ بالنسبة إلى $(D)$.
\QQ{2} $g$ الدالة المعرفة على $\R$ بـ: $g(x)=(2x-1)e^{2x}-1$. استنتج حسب قيم $x$ إشارة $g(x)$، ثم تحقّق أن: $0{,}6<\alpha<0{,}7$.

\textbf{II)} $f$ الدالة المعرفة على $\R$ بـ: $f(x)=(x-1)\left(e^{2x}-1\right)$، و $(C_f)$ تمثيلها البياني في المستوي المنسوب إلى معلم متعامد ومتجانس $(O;\vec i,\vec j)$ (الوحدة $2cm$).
\QQ{1} أحسب $\lim\limits_{x\to+\infty}f(x)$ و $\lim\limits_{x\to-\infty}f(x)$.
\QQ{2}
\begin{sq}
\qa بيّن أن المستقيم $(\Delta)$ ذا المعادلة $y=-x+1$ مقارب مائل لـ $(C_f)$ عند $-\infty$.
\qb أدرس وضعية $(C_f)$ بالنسبة إلى $(\Delta)$.
\end{sq}
\QQ{3}
\begin{sq}
\qa بيّن أنه من أجل كل عدد حقيقي $x$: $f'(x)=g(x)$.
\qb استنتج أن الدالة $f$ متناقصة تماما على المجال $]{-}\infty;\alpha]$ ومتزايدة تماما على المجال $[\alpha;+\infty[$ ثم شكّل جدول تغيّراتها.
\qc بيّن أن $(C_f)$ يقبل مماسا $(T)$ موازيا لـ $(\Delta)$، يُطلب تعيين معادلة له.
\end{sq}
\QQ{4}
\begin{sq}
\qa عيّن فواصل نقط تقاطع $(C_f)$ مع حامل محور الفواصل.
\qb أرسم $(\Delta)$، $(T)$ و $(C_f)$. (نأخذ: $f(1{,}4)\approx6{,}2$ و $f(\alpha)\approx-0{,}9$).
\qc ناقش بيانيا، حسب قيم الوسيط الحقيقي $m$، عدد حلول المعادلة $f(x)=-x+m$.
\end{sq}
\end{tamrin}

% ===================== بكالوريا 2022 -- الموضوع الثاني =====================
\begin{tamrin}{بكالوريا 2022 -- الموضوع الثاني -- علوم تجريبية}
$f$ الدالة العددية المعرفة على $\R$ بـ: $f(x)=\dfrac12e^{-2x}-\dfrac92e^{-x}-2x+4$، و $(C_f)$ تمثيلها البياني في مستوٍ منسوب إلى المعلم المتعامد والمتجانس $(O;\vec i,\vec j)$.
\QQ{1} أحسب $\lim\limits_{x\to+\infty}f(x)$ وبيّن أن: $\lim\limits_{x\to-\infty}f(x)=+\infty$.
\QQ{2}
\begin{sq}
\qa أثبت أنه من أجل كل عدد حقيقي $x$: $f'(x)=-\dfrac12e^{-2x}\left(e^{x}-2\right)\left(4e^{x}-1\right)$.
\qb بيّن أن $f$ متناقصة على كل من المجالين $]{-}\infty;-\ln4]$ و $[\ln2;+\infty[$ ومتزايدة على $[-\ln4;\ln2]$، ثم شكّل جدول تغيّراتها.
\end{sq}
\QQ{3}
\begin{sq}
\qa بيّن أن المستقيم $(\Delta)$ ذا المعادلة $y=-2x+4$ مقارب مائل للمنحنى $(C_f)$ عند $+\infty$.
\qb أدرس وضعية $(C_f)$ بالنسبة إلى $(\Delta)$.
\end{sq}
\QQ{4} أكتب معادلة لـ $(T)$ مماس $(C_f)$ في النقطة ذات الفاصلة $0$.
\QQ{5} أنشئ $(\Delta)$ و $(T)$ والمنحنى $(C_f)$ على المجال $[-1{,}9;+\infty[$. (نأخذ $f(-1{,}9)\approx0$ و $f(-\ln4)\approx-3{,}2$).
\QQ{6} $h$ الدالة المعرفة على $\R$ بـ: $h(x)=-\dfrac12e^{-2x}+\dfrac92e^{-x}+2x-2$، و $(C_h)$ تمثيلها البياني في المعلم السابق.
\begin{sq}
\qa عيّن العددين الحقيقيين $a$ و $b$ حيث، من أجل كل عدد حقيقي $x$: $h(x)=af(x)+b$.
\qb اشرح كيف يمكن إنشاء $(C_h)$ اعتمادا على $(C_f)$ (لا يُطلب إنشاء $(C_h)$).
\end{sq}
\end{tamrin}

% ===================== بكالوريا 2022 -- الموضوع الأول =====================
\begin{tamrin}{بكالوريا 2022 -- الموضوع الأول -- علوم تجريبية}
$f$ الدالة العددية المعرفة على $\R$ بتمثيلها البياني $(C_f)$ في المستوي المنسوب إلى المعلم المتعامد والمتجانس $(O;\vec i,\vec j)$، و $(T)$ مماس $(C_f)$ في النقطة ذات الفاصلة $0$ كما هو مبيّن في الشكل المقابل.
\begin{GR}{-5.5}{1.5}{-1.6}{3.5}{0.8}
\Ticks{-5,-4,-3,-2,-1,1}{-1,1,2,3}
\draw[cRem,thick](-5.5,-1)--(1.5,-1);
\draw[cRem!60!black,thick,domain=-1.6:1.5] plot(\x,{\x});
\node[cRem!60!black] at (1.1,0.5){$(T)$};
\draw[cDef,very thick,domain=-5.5:0.95,samples=150] plot(\x,{(\x*\x+1)*exp(\x)-1});
\node[cDef] at (0.3,2.8){$(C_f)$};
\end{GR}
\QQ{1} بقراءة بيانية: عيّن $f'(0)$ و $\lim\limits_{x\to-\infty}f(x)$، وأعط معادلة للمماس $(T)$.
\QQ{2} ناقش بيانيا، حسب قيم الوسيط الحقيقي $m$، عدد حلول المعادلة $f(x)=x+m$.
\QQ{3} بيّن أن: $a=1$ و $b=-1$ إذا علمت أن: $f(x)=\left(x^{2}+a\right)e^{x}+b$.
\QQ{4} $g$ الدالة العددية المعرفة على $\R$ بـ: $g(x)=\left(x^{2}+1\right)e^{|x|}-1$، و $(C_g)$ تمثيلها البياني في المعلم السابق.
بيّن أن الدالة $g$ زوجية، ثم اشرح كيفية إنشاء $(C_g)$ انطلاقا من $(C_f)$ ثم أنشئ $(C_g)$.
\end{tamrin}

% ===================== بكالوريا 2021 -- الموضوع الثاني =====================
\begin{tamrin}{بكالوريا 2021 -- الموضوع الثاني -- علوم تجريبية}
\textbf{I)} الدالة العددية $g$ معرفة على $\R$ بـ: $g(x)=1+xe^{-x-1}$، و $(C_g)$ تمثيلها البياني في المستوي المنسوب إلى المعلم المتعامد والمتجانس $(O;\vec i,\vec j)$ (الشكل).
\begin{GR}{-2.5}{3.8}{-1.5}{1.6}{1}
\Ticks{-2,-1,1,2,3}{-1,1}
\draw[cDef,very thick,domain=-1.8:3.8,samples=150] plot(\x,{1+\x*exp(-\x-1)});
\node[cDef] at (-1.5,0.9){$(C_g)$};
\end{GR}
\QQ{1} أحسب $g(-1)$.
\QQ{2} بقراءة بيانية، حدّد حسب قيم $x$ إشارة $g(x)$.

\textbf{II)} الدالة العددية $f$ معرفة على $\R$ بـ: $f(x)=x-(x+1)e^{-x-1}$، و $(C_f)$ تمثيلها البياني في المستوي المنسوب إلى المعلم المتعامد والمتجانس $(O;\vec i,\vec j)$.
\QQ{1} تحقّق أنه من أجل كل عدد حقيقي $x$ غير معدوم: $f(x)=x\left[1-\left(1+\dfrac1x\right)e^{-x-1}\right]$، ثم أحسب $\lim\limits_{x\to+\infty}f(x)$ و $\lim\limits_{x\to-\infty}f(x)$.
\QQ{2}
\begin{sq}
\qa بيّن أنه من أجل كل عدد حقيقي $x$: $f'(x)=g(x)$.
\qb استنتج أن الدالة $f$ متزايدة تماما على المجال $[-1;+\infty[$ ومتناقصة تماما على المجال $]{-}\infty;-1]$ ثم شكّل جدول تغيّراتها.
\end{sq}
\QQ{3}
\begin{sq}
\qa أحسب $\lim\limits_{x\to+\infty}\left[f(x)-x\right]$، ثم فسّر النتيجة هندسيا.
\qb أدرس وضعية $(C_f)$ بالنسبة للمستقيم $(\Delta)$ ذي المعادلة $y=x$.
\qc بيّن أن المنحنى $(C_f)$ يقبل مماسا $(T)$ موازيا للمستقيم $(\Delta)$ يُطلب كتابة معادلة له.
\end{sq}
\QQ{4}
\begin{sq}
\qa بيّن أن $(C_f)$ يقطع حامل محور الفواصل في نقطتين فاصلتاهما $\alpha$ و $\beta$ حيث: $0{,}3<\alpha<0{,}4$ و $-1{,}9<\beta<-1{,}8$.
\qb أرسم المستقيمين $(\Delta)$ و $(T)$ ثم أرسم المنحنى $(C_f)$ على المجال $]-2;+\infty[$.
\end{sq}
\QQ{5} الدالة العددية $h$ معرفة على المجال $[-2;2]$ بـ: $h(x)=-|x|+\left(|x|-1\right)e^{|x|-1}$، و $(C_h)$ تمثيلها البياني في المعلم السابق.
\begin{sq}
\qa أثبت أن الدالة $h$ زوجية.
\qb بيّن أنه من أجل كل $x\in[-2;0]$: $h(x)=f(x)$.
\qc اشرح كيفية رسم $(C_h)$ انطلاقا من $(C_f)$، ثم ارسمه.
\end{sq}
\end{tamrin}

\begin{molahada}[إتمام نص التمرين]
في نص المجلة ينتهي السؤال 5 بتعريف الدالة $h$ دون أي مطلوب؛ أُتمّ بالأسئلة (أ)، (ب)، (جـ) أعلاه كما وردت في حل التمرين.
\end{molahada}

% ===================== بكالوريا 2020 -- الموضوع الثاني =====================
\begin{tamrin}{بكالوريا 2020 -- الموضوع الثاني -- علوم تجريبية}
المستوي منسوب إلى المعلم المتعامد والمتجانس $(O;\vec i,\vec j)$.

\textbf{I)} في الشكل المرفق، $(\Gamma)$ المنحنى الممثل للدالة $g$ المعرفة على $\R$ بـ: $g(x)=2x^{2}+2x-2xe^{x}$، و $(\Delta)$ المستقيم ذو المعادلة $y=x$، و $(\gamma)$ المنحنى الممثل للدالة $x\mapsto e^{x}$.
\begin{GR}{-3}{2}{-2}{3}{0.9}
\Ticks{-2,-1,1}{-1,1,2}
\draw[cThm,thick,domain=-2:2] plot(\x,{\x}) node[pos=0.06,above left]{$(\Delta)$};
\draw[cRem!70!black,thick,domain=-3:1.1,samples=100] plot(\x,{exp(\x)});
\node[cRem!70!black] at (1.5,2.6){$(\gamma)$};
\draw[cDef,very thick,domain=-3:1.2,samples=150] plot(\x,{2*\x*\x+2*\x-2*\x*exp(\x)});
\node[cDef] at (-2.3,2.6){$(\Gamma)$};
\end{GR}
بقراءة بيانية:
\QQ{1} برّر أنه من أجل كل عدد حقيقي $x$: $e^{x}-x>0$.
\QQ{2} حدّد تبعا لقيم العدد الحقيقي $x$ إشارة $g(x)$ علما أن: $g(0)=0$.

\textbf{II)} $f$ الدالة المعرفة على $\R$ كما يلي: $f(x)=-1+\dfrac{2e^{x}}{e^{x}-x}$، و $(C_f)$ تمثيلها البياني في المستوي السابق.
\QQ{1} بيّن أن $\lim\limits_{x\to+\infty}f(x)=1$ وأحسب $\lim\limits_{x\to-\infty}f(x)$، ثم فسّر النتيجتين هندسيا.
\QQ{2}
\begin{sq}
\qa بيّن أنه من أجل كل عدد حقيقي $x$: $f'(x)=\dfrac{2e^{x}(1-x)}{\left(e^{x}-x\right)^{2}}$.
\qb استنتج اتجاه تغيّر الدالة $f$، ثم شكّل جدول تغيّراتها.
\end{sq}
\QQ{3}
\begin{sq}
\qa أكتب معادلة لـ $(T)$ مماس المنحنى $(C_f)$ في النقطة $A$ ذات الفاصلة $0$.
\qb بيّن أنه من أجل كل عدد حقيقي $x$: $f(x)-(2x+1)=\dfrac{g(x)}{e^{x}-x}$.
\qc استنتج الوضع النسبي لـ $(C_f)$ و $(T)$، ماذا تمثل النقطة $A$ بالنسبة إلى $(C_f)$؟
\end{sq}
\QQ{4} بيّن أن المعادلة $f(x)=0$ تقبل حلا وحيدا $\alpha$ في المجال $]{-}\infty;1[$، ثم تحقّق أن: $-0{,}6<\alpha<-0{,}5$.
\QQ{5} أنشئ المماس $(T)$ والمستقيمين المقاربين ثم المنحنى $(C_f)$.
\end{tamrin}

% ===================== بكالوريا 2019 -- الموضوع الثاني =====================
\begin{tamrin}{بكالوريا 2019 -- الموضوع الثاني -- علوم تجريبية}
المستوي منسوب إلى المعلم المتعامد والمتجانس $(O;\vec i,\vec j)$، وتؤخذ وحدة الطول $2cm$.
$(C_f)$ و $(C_g)$ التمثيلان البيانيان للدالتين $f$ و $g$ المعرفتين على $\R$ كما يلي:
\[f(x)=e^{x}-\frac12ex^{2}\qquad;\qquad g(x)=e^{x}-ex\]
\QQ{1}
\begin{sq}
\qa أدرس اتجاه تغيّر الدالة $g$.
\qb استنتج إشارة $g(x)$ حسب قيم $x$.
\end{sq}
\QQ{2} أدرس اتجاه تغيّر الدالة $f$.
\QQ{3} أحسب كلا من $\lim\limits_{x\to-\infty}f(x)$ و $\lim\limits_{x\to+\infty}f(x)$، ثم شكّل جدول تغيّرات الدالة $f$.
\QQ{4} أدرس الوضع النسبي للمنحنيين $(C_f)$ و $(C_g)$ على $\R$.
\QQ{5} أرسم على المجال $[0;2]$ المنحنيين $(C_f)$ و $(C_g)$ في نفس المعلم $(O;\vec i,\vec j)$. (يُعطى $e^{2}-2e\approx2$).
\QQ{6} $h$ الدالة المعرفة على المجال $[-2;2]$ كما يلي: $h(x)=\dfrac12ex^{2}-e^{|x|}$، وليكن $(\Gamma)$ تمثيلها البياني في المعلم السابق.
\begin{sq}
\qa بيّن أن الدالة $h$ زوجية.
\qb من أجل $x\in[0;2]$ أحسب $h(x)+f(x)$، ثم استنتج كيفية رسم $(\Gamma)$ انطلاقا من $(C_f)$ ثم أرسمه.
\end{sq}
\end{tamrin}

% ===================== بكالوريا 2018 -- الموضوع الأول =====================
\begin{tamrin}{بكالوريا 2018 -- الموضوع الأول -- علوم تجريبية}
\textbf{I)} $g$ الدالة العددية المعرفة على $\R$ كما يلي: $g(x)=2+(x-1)e^{-x}$.
\QQ{1} أحسب $\lim\limits_{x\to-\infty}g(x)$ و $\lim\limits_{x\to+\infty}g(x)$.
\QQ{2} أدرس اتجاه تغيّر الدالة $g$ ثم شكّل جدول تغيّراتها.
\QQ{3} بيّن أن المعادلة $g(x)=0$ تقبل حلا وحيدا $\alpha$ حيث $-0{,}38<\alpha<-0{,}37$، ثم استنتج إشارة $g(x)$ على $\R$.

\textbf{II)} نعتبر الدالة $f$ المعرفة على $\R$ بـ: $f(x)=2x+1-xe^{-x}$، و $(C_f)$ تمثيلها البياني في المستوي المنسوب إلى المعلم المتعامد والمتجانس $(O;\vec i,\vec j)$.
\QQ{1}
\begin{sq}
\qa أحسب $\lim\limits_{x\to-\infty}f(x)$ و $\lim\limits_{x\to+\infty}f(x)$.
\qb أحسب $\lim\limits_{x\to+\infty}\left[f(x)-(2x+1)\right]$، ثم فسّر النتيجة بيانيا.
\qc أدرس الوضع النسبي للمنحنى $(C_f)$ والمستقيم $(\Delta)$ حيث $(\Delta):y=2x+1$.
\end{sq}
\QQ{2} بيّن أنه من أجل كل عدد حقيقي $x$ يكون: $f'(x)=g(x)$، ثم استنتج اتجاه تغيّر الدالة $f$ وشكّل جدول تغيّراتها.
\QQ{3} أكتب معادلة المماس $(T)$ للمنحنى $(C_f)$ عند النقطة ذات الفاصلة $1$.
\QQ{4} أنشئ $(\Delta)$، $(T)$ والمنحنى $(C_f)$. (نأخذ $f(\alpha)\approx0{,}8$).
\QQ{5} ناقش بيانيا وحسب قيم الوسيط الحقيقي $m$ عدد وإشارة حلول المعادلة ذات المجهول $x$: $x=(1-m)e^{x}$.
\end{tamrin}

% ===================== بكالوريا 2017 -- الموضوع الأول -- الدورة الاستثنائية =====================
\begin{tamrin}{بكالوريا 2017 -- الموضوع الأول -- علوم تجريبية -- الدورة الاستثنائية}
\textbf{I)} نعتبر الدالة $g$ المعرفة على $\R$ بـ: $g(x)=x^{2}e^{x}-e$، و $(C_g)$ تمثيلها البياني في المستوي المنسوب إلى المعلم المتعامد والمتجانس $(O;\vec i,\vec j)$ (الشكل).
\begin{GR}{-6}{1.6}{-3.2}{1.5}{0.85}
\Ticks{-5,-4,-3,-2,-1,1}{-3,-2,-1,1}
\draw[cDef,very thick,domain=-6:1.25,samples=200] plot(\x,{\x*\x*exp(\x)-exp(1)});
\node[cDef] at (-4.8,-1.6){$(C_g)$};
\end{GR}
\QQ{1} أحسب $g(1)$.
\QQ{2} بقراءة بيانية، عيّن إشارة $g(x)$، ثم استنتج إشارة $g(-x)$ حسب قيم العدد الحقيقي $x$.

\textbf{II)} نعتبر الدالة $f$ المعرفة على $\R^{*}$ بـ: $f(x)=e^{-x}-2-\dfrac{e}{x}$، و $(C_f)$ تمثيلها البياني في المستوي المنسوب إلى المعلم المتعامد والمتجانس $(O;\vec i,\vec j)$.
\QQ{1} أحسب النهايات الآتية: $\lim\limits_{x\to-\infty}f(x)$، $\lim\limits_{x\to0^{-}}f(x)$، $\lim\limits_{x\to0^{+}}f(x)$ و $\lim\limits_{x\to+\infty}f(x)$.
\QQ{2} بيّن أن المنحنى $(\gamma)$ الذي معادلته $y=e^{-x}-2$ والمنحنى $(C_f)$ متقاربان بجوار $+\infty$، ثم أدرس وضعية المنحنى $(C_f)$ بالنسبة لـ $(\gamma)$.
\QQ{3} بيّن أنه من أجل كل $x$ من $\R^{*}$: $f'(x)=\dfrac{-g(-x)}{x^{2}}$.
\QQ{4} استنتج أن الدالة $f$ متزايدة تماما على كل من المجالين $]-1;0[$ و $]0;+\infty[$ ومتناقصة تماما على المجال $]{-}\infty;-1]$، ثم شكّل جدول تغيّراتها.
\QQ{5} بيّن كيف يمكن إنشاء المنحنى $(\gamma)$ انطلاقا من منحنى الدالة $x\mapsto e^{x}$، ثم أرسم بعناية كلا من $(\gamma)$ و $(C_f)$ في نفس المعلم السابق.
\end{tamrin}

% ===================== بكالوريا 2017 م2 =====================
\begin{tamrin}{بكالوريا 2017 -- الموضوع الثاني -- علوم تجريبية}
\textbf{I)} نعتبر الدالة العددية $f$ المعرفة على $\R$ بـ: $f(x)=2-x^2e^{1-x}$، و $(C_f)$ تمثيلها البياني في المستوي المنسوب إلى المعلم المتعامد والمتجانس $(O;\vec i,\vec j)$.

\QQ{1} بيّن أن $\lim\limits_{x\to+\infty}f(x)=2$ وأعط تفسيرًا هندسيًا للنتيجة، ثم أحسب $\lim\limits_{x\to-\infty}f(x)$.
\QQ{2}
\begin{sq}
\qa بيّن أنه من أجل كل $x$ من $\R$: $f'(x)=x(x-2)e^{1-x}$.
\qb أدرس اتجاه تغير الدالة $f$ ثم شكّل جدول تغيراتها.
\end{sq}
\QQ{3} أكتب معادلة لـ $(T)$ المماس للمنحنى $(C_f)$ عند النقطة ذات الفاصلة $1$.

\textbf{II)} $h$ الدالة العددية المعرفة على $\R$ بـ: $h(x)=1-xe^{1-x}$.

\QQ{1} بيّن أنه من أجل كل $x$ من $\R$: $h(x)\geqslant 0$، ثم أدرس الوضع النسبي للمنحنى $(C_f)$ والمماس $(T)$.
\QQ{2} بيّن أن المعادلة $f(x)=0$ تقبل حلًا وحيدًا $\alpha$ حيث $-0{,}7<\alpha<-0{,}6$.
\QQ{3} أنشئ المماس $(T)$ والمنحنى $(C_f)$ على المجال $[-1;+\infty[$.
\QQ{4} $F$ الدالة المعرفة على $\R$ بـ: $F(x)=2x+(x^2+2x+2)e^{1-x}$. بيّن أنه من أجل كل عدد حقيقي $x$: $F'(x)=f(x)$.
\end{tamrin}

% ===================== بكالوريا 2016 م2 د2 =====================
\begin{tamrin}{بكالوريا 2016 -- الموضوع الثاني -- علوم تجريبية -- الدورة الثانية}
\textbf{I)} $g$ الدالة العددية المعرفة على $\R$ بـ: $g(x)=2e^x-x^2-x$.

\QQ{1}
\begin{sq}
\qa أحسب $g'(x)$ من أجل كل $x$ من $\R$، ثم أدرس اتجاه تغير الدالة $g'$ (حيث $g'$ هي مشتقة الدالة $g$).
\qb بيّن أنه من أجل كل $x$ من $\R$: $g'(x)>0$.
\qc أحسب نهايتي الدالة $g$ عند كل من $+\infty$ و $-\infty$، ثم شكّل جدول تغيراتها.
\end{sq}
\QQ{2} بيّن أن المعادلة $g(x)=0$ تقبل حلًا وحيدًا $\alpha$ حيث: $-1{,}38<\alpha<-1{,}37$.
\QQ{3} استنتج إشارة $g(x)$ حسب قيم العدد الحقيقي $x$.

\textbf{II)} $f$ الدالة العددية المعرفة على $\R$ كما يلي: $f(x)=\dfrac{x^2e^x}{e^x-x}$، و $(C_f)$ تمثيلها البياني في المستوي المنسوب إلى المعلم المتعامد والمتجانس $(O;\vec i,\vec j)$.

\QQ{1}
\begin{sq}
\qa أحسب $\lim\limits_{x\to+\infty}f(x)$ و $\lim\limits_{x\to-\infty}f(x)$.
\qb بيّن أنه من أجل كل $x$ من $\R$: $f'(x)=\dfrac{xe^xg(x)}{(e^x-x)^2}$.
\qc أدرس اتجاه تغير الدالة $f$ على $\R$، ثم شكّل جدول تغيراتها.
\end{sq}
\QQ{2}
\begin{sq}
\qa بيّن أن: $f(\alpha)=\alpha^2+2\alpha+2+\dfrac{2}{\alpha-1}$، ثم استنتج حصرًا للعدد $f(\alpha)$.
\qb أحسب $\lim\limits_{x\to+\infty}\left[f(x)-x^2\right]$، ثم فسّر النتيجة بيانيًا.
\qc أنشئ المنحنى $(C_f)$. (تعطى $f(\alpha)\approx0{,}29$)
\end{sq}
\end{tamrin}

% ===================== بكالوريا 2016 م2 د1 =====================
\begin{tamrin}{بكالوريا 2016 -- الموضوع الثاني -- علوم تجريبية -- الدورة الأولى}
\textbf{I)} لتكن $g$ الدالة العددية المعرفة على $\R$ بـ: $g(x)=1+(x^2+x-1)e^{-x}$.

\QQ{1} أحسب $\lim\limits_{x\to-\infty}g(x)$ و $\lim\limits_{x\to+\infty}g(x)$.
\QQ{2} أدرس اتجاه تغير الدالة $g$، ثم شكّل جدول تغيراتها.
\QQ{3}
\begin{sq}
\qa بيّن أن للمعادلة $g(x)=0$ حلّين في $\R$، أحدهما معدوم والآخر $\alpha$ حيث: $-1{,}52<\alpha<-1{,}51$.
\qb استنتج إشارة $g(x)$ على $\R$.
\end{sq}

\textbf{II)} نعتبر الدالة $f$ المعرفة على $\R$ كما يلي: $f(x)=-x+(x^2+3x+2)e^{-x}$، و $(C_f)$ تمثيلها البياني في المستوي المنسوب إلى المعلم المتعامد والمتجانس $(O;\vec i,\vec j)$.

\QQ{1}
\begin{sq}
\qa أحسب $\lim\limits_{x\to-\infty}f(x)$ و $\lim\limits_{x\to+\infty}f(x)$.
\qb بيّن أنه من أجل كل عدد حقيقي $x$ فإن: $f'(x)=-g(x)$.
\qc شكّل جدول تغيرات الدالة $f$ على $\R$. (نأخذ $f(\alpha)\approx0{,}38$)
\qd عيّن دون حساب: $\lim\limits_{h\to0}\dfrac{f(\alpha+h)-f(\alpha)}{h}$، ثم فسّر النتيجة هندسيًا.
\end{sq}
\QQ{2}
\begin{sq}
\qa بيّن أن المستقيم $(\Delta)$ ذا المعادلة $y=-x$ مقارب مائل للمنحنى $(C_f)$ عند $+\infty$.
\qb أدرس وضعية المنحنى $(C_f)$ بالنسبة للمستقيم $(\Delta)$.
\qc بيّن أن المنحنى $(C_f)$ يقبل نقطتي انعطاف يطلب تعيين إحداثييهما.
\qd أرسم $(\Delta)$ و $(C_f)$ على المجال $[-2;+\infty[$.
\qe ناقش بيانيًا وحسب قيم الوسيط $m$ عدد وإشارة حلول المعادلة: $(m-x)e^x+(x^2+3x+2)=0$ على المجال $[-2;+\infty[$.
\end{sq}
\end{tamrin}

% ===================== بكالوريا 2015 م2 =====================
\begin{tamrin}{بكالوريا 2015 -- الموضوع الثاني -- علوم تجريبية}
\textbf{I)} $g$ الدالة العددية المعرفة على $\R$ بـ: $g(x)=1-2x-e^{2x-2}$.

\QQ{1} أدرس اتجاه تغير الدالة $g$ على $\R$.
\QQ{2} بيّن أن المعادلة $g(x)=0$ تقبل حلًا وحيدًا $\alpha$ في $\R$، ثم تحقق أن: $0{,}36<\alpha<0{,}37$.
\QQ{3} استنتج إشارة $g(x)$ على $\R$.

\textbf{II)} $f$ الدالة العددية المعرفة على $\R$ كما يلي: $f(x)=xe^{2x+2}-x+1$، و $(C_f)$ تمثيلها البياني في المستوي المنسوب إلى المعلم المتعامد والمتجانس $(O;\vec i,\vec j)$.

\QQ{1}
\begin{sq}
\qa بيّن أنه من أجل كل عدد حقيقي $x$: $f'(x)=e^{2x+2}g(-x)$.
\qb استنتج أن الدالة $f$ متناقصة تمامًا على المجال $]{-}\infty;-\alpha]$ ومتزايدة تمامًا على $[-\alpha;+\infty[$.
\end{sq}
\QQ{2} أحسب نهاية $f$ عند $+\infty$ وعند $-\infty$، ثم شكّل جدول تغيرات الدالة $f$.
\QQ{3} أحسب $\lim\limits_{x\to-\infty}\left[f(x)+x-1\right]$، ثم فسّر النتيجة بيانيًا.
\QQ{4} أدرس وضعية المنحنى $(C_f)$ بالنسبة إلى المستقيم $(\Delta)$ الذي معادلته $y=-x+1$.
\QQ{5} أنشئ $(\Delta)$ و $(C_f)$ على المجال $\left]{-}\infty;\frac12\right]$. (نأخذ $f(-\alpha)\approx0{,}1$)
\QQ{6} تحقق أنه من أجل كل $x$ من $\R$: $2f(x)+f'(x)-f''(x)=1-2x-3e^{2x+2}$.
\end{tamrin}

% ===================== بكالوريا 2013 م1 =====================
\begin{tamrin}{بكالوريا 2013 -- الموضوع الأول -- علوم تجريبية}
\textbf{I)} $f$ الدالة المعرفة على المجال $]{-}\infty;1[$ بـ: $f(x)=\dfrac{x}{x-1}+e^{\frac{1}{x-1}}$، و $(C)$ تمثيلها البياني في المستوي المنسوب إلى المعلم المتعامد والمتجانس $(O;\vec i,\vec j)$.

\QQ{1} أحسب $\lim\limits_{x\to-\infty}f(x)$ و $\lim\limits_{\substack{x\to1\\x<1}}f(x)$، ثم استنتج المستقيمين المقاربين للمنحنى $(C)$.
\QQ{2} أحسب $f'(x)$. بيّن أن الدالة $f$ متناقصة تمامًا على المجال $]{-}\infty;1[$، ثم شكّل جدول تغيراتها.
\QQ{3} بيّن أن المعادلة $f(x)=0$ تقبل في $]{-}\infty;1[$ حلًا وحيدًا $\alpha$. باستعمال جدول القيم الآتي جد حصرًا للعدد $\alpha$:
\begin{center}
\begin{tabular}{|c|c|c|c|c|c|c|}\hline
\rowcolor{cSec!12}$x$ & $0{,}20$ & $0{,}21$ & $0{,}22$ & $0{,}23$ & $0{,}24$ & $0{,}25$\\\hline
$f(x)$ & $0{,}037$ & $0{,}016$ & $-0{,}005$ & $-0{,}026$ & $-0{,}048$ & $-0{,}070$\\\hline
\end{tabular}
\end{center}
\QQ{4} أرسم المستقيمين المقاربين والمنحنى $(C)$، ثم أرسم المنحنى $(C')$ الممثل للدالة $|f|$.
\QQ{5} عيّن بيانيًا مجموعة قيم الأعداد الحقيقية $m$ التي من أجلها يكون للمعادلة $|f(x)|=m$ حلّان مختلفان في الإشارة.

\textbf{II)} $g$ الدالة المعرفة على المجال $]{-}\infty;1[$ بـ: $g(x)=f(2x-1)$. (عبارة $g(x)$ غير مطلوبة)

\QQ{1} أدرس تغيرات الدالة $g$ على $]{-}\infty;1[$، ثم شكّل جدول تغيراتها.
\QQ{2}
\begin{sq}
\qa تحقق أن $g\left(\dfrac{\alpha+1}{2}\right)=0$، ثم بيّن أن: $g'\left(\dfrac{\alpha+1}{2}\right)=2f'(\alpha)$.
\qb استنتج معادلة $(T)$ المماس لمنحنى الدالة $g$ في النقطة ذات الفاصلة $\dfrac{\alpha+1}{2}$.
\qc تحقق من أن: $y=\dfrac{2}{(\alpha-1)^3}x-\dfrac{\alpha+1}{(\alpha-1)^3}$ معادلة للمستقيم $(T)$.
\end{sq}
\end{tamrin}

% ===================== بكالوريا 2012 م2 =====================
\begin{tamrin}{بكالوريا 2012 -- الموضوع الثاني -- علوم تجريبية}
\textbf{I)} لتكن $g$ الدالة العددية المعرفة على $\R$ بـ: $g(x)=1-xe^x$.

\QQ{1} أحسب $\lim\limits_{x\to-\infty}g(x)$ و $\lim\limits_{x\to+\infty}g(x)$.
\QQ{2} أدرس اتجاه تغير الدالة $g$ ثم شكّل جدول تغيراتها.
\QQ{3}
\begin{sq}
\qa بيّن أن المعادلة $g(x)=0$ تقبل حلًا وحيدًا $\alpha$ على المجال $]-1;+\infty[$.
\qb تحقق أن $0{,}5<\alpha<0{,}6$، ثم استنتج إشارة $g(x)$ على $\R$.
\end{sq}

\textbf{II)} نعتبر الدالة $f$ المعرفة على المجال $]{-}\infty;2]$ كما يلي: $f(x)=(x-1)e^x-x-1$، و $(C_f)$ تمثيلها البياني في المستوي المنسوب إلى المعلم المتعامد والمتجانس $(O;\vec i,\vec j)$.

\QQ{1} أحسب $\lim\limits_{x\to-\infty}f(x)$.
\QQ{2} لتكن $f'$ مشتقة الدالة $f$. بيّن أنه من أجل كل عدد حقيقي $x$ من $]{-}\infty;2]$ فإن: $f'(x)=-g(x)$، استنتج إشارة $f'(x)$ على المجال $]{-}\infty;2]$، ثم شكّل جدول تغيرات الدالة $f$.
\QQ{3} بيّن أن $f(\alpha)=-\left(\dfrac{\alpha^2+1}{\alpha}\right)$، ثم استنتج حصرًا للعدد $f(\alpha)$. (تدوّر النتائج إلى $10^{-2}$)
\QQ{4}
\begin{sq}
\qa بيّن أن المستقيم $(\Delta)$ ذا المعادلة $y=-x-1$ هو مستقيم مقارب مائل للمنحنى $(C_f)$ عند $-\infty$.
\qb أدرس وضعية $(C_f)$ بالنسبة إلى $(\Delta)$.
\end{sq}
\QQ{5}
\begin{sq}
\qa بيّن أن المعادلة $f(x)=0$ تقبل حلّين $x_1$ و $x_2$ حيث $-1{,}6<x_1<-1{,}5$ و $1{,}5<x_2<1{,}6$.
\qb أنشئ $(\Delta)$ و $(C_f)$.
\end{sq}
\end{tamrin}

% ===================== بكالوريا 2011 م2 =====================
\begin{tamrin}{بكالوريا 2011 -- الموضوع الثاني -- علوم تجريبية}
نعتبر الدالة العددية $f$ المعرفة على $\R$ بـ: $f(x)=e^x-ex-1$، و $(C_f)$ تمثيلها البياني في المستوي المنسوب إلى المعلم المتعامد والمتجانس $(O;\vec i,\vec j)$.

\QQ{1}
\begin{sq}
\qa أحسب $\lim\limits_{x\to-\infty}f(x)$ و $\lim\limits_{x\to+\infty}f(x)$.
\qb أحسب $f'(x)$ ثم أدرس إشارتها.
\qc شكّل جدول تغيرات الدالة $f$.
\end{sq}
\QQ{2}
\begin{sq}
\qa بيّن أن المستقيم $(\Delta)$ ذا المعادلة $y=-ex-1$ مقارب مائل للمنحنى $(C_f)$ عند $-\infty$.
\qb أكتب معادلة للمستقيم $(T)$ مماس المنحنى $(C_f)$ عند النقطة ذات الفاصلة $0$.
\qc بيّن أن المعادلة $f(x)=0$ تقبل في المجال $]1{,}75;1{,}76[$ حلًا وحيدًا $\alpha$.
\qd أرسم المستقيمين $(\Delta)$ و $(T)$ ثم المنحنى $(C_f)$ على المجال $]{-}\infty;2]$.
\end{sq}
\end{tamrin}

% ===================== بكالوريا 2010 م2 =====================
\begin{tamrin}{بكالوريا 2010 -- الموضوع الثاني -- علوم تجريبية}
نعتبر الدالة العددية $f$ المعرفة على $\R^*$ كما يلي: $f(x)=x-\dfrac{1}{e^x-1}$.

نرمز بـ $(C_f)$ لتمثيلها البياني في المستوي المنسوب إلى المعلم المتعامد والمتجانس $(O;\vec i,\vec j)$.

\QQ{1}
\begin{sq}
\qa أحسب $\lim\limits_{x\to-\infty}f(x)$ و $\lim\limits_{x\to+\infty}f(x)$.
\qb أحسب $\lim\limits_{\substack{x\to0\\x>0}}f(x)$ و $\lim\limits_{\substack{x\to0\\x<0}}f(x)$ وفسّر هندسيًا النتيجة.
\end{sq}
\QQ{2} أدرس اتجاه تغير الدالة $f$ على كل مجال من مجالي تعريفها ثم شكّل جدول تغيراتها.
\QQ{3}
\begin{sq}
\qa بيّن أن المنحنى $(C_f)$ يقبل مستقيمين مقاربين مائلين $(\Delta)$ و $(\Delta')$ معادلتاهما على الترتيب: $y=x$ و $y=x+1$.
\qb أدرس وضعية $(C_f)$ بالنسبة إلى كل من $(\Delta)$ و $(\Delta')$.
\end{sq}
\QQ{4} أثبت أن النقطة $\omega\left(0;\frac12\right)$ هي مركز تناظر للمنحنى $(C_f)$.
\QQ{5}
\begin{sq}
\qa بيّن أن المعادلة $f(x)=0$ تقبل حلّين $\alpha$ و $\beta$ حيث: $\ln2<\alpha<1$ و $-1{,}4<\beta<-1{,}3$.
\qb هل توجد مماسات لـ $(C_f)$ توازي المستقيم $(\Delta)$؟
\qc أرسم $(\Delta)$، $(\Delta')$ ثم المنحنى $(C_f)$.
\qd ناقش بيانيًا حسب قيم الوسيط $m$ عدد وإشارة حلول المعادلة: $(m-1)e^{-x}=m$.
\end{sq}
\end{tamrin}

% ===================== بكالوريا 2008 م1 =====================
\begin{tamrin}{بكالوريا 2008 -- الموضوع الأول -- علوم تجريبية}
\textbf{I)} نعتبر الدالة العددية $f$ للمتغير الحقيقي $x$ المعرفة على المجال $]-2;+\infty[$ كما يلي: $f(x)=(ax+b)e^{-x}+1$.

$(C_f)$ تمثيلها البياني في المستوي المنسوب إلى المعلم المتعامد والمتجانس $(O;\vec i,\vec j)$ (وحدة الطول $1\,cm$).

عيّن قيمتي $a$ و $b$ بحيث تكون النقطة $A(-1;1)$ تنتمي إلى $(C_f)$ ومعامل توجيه المماس عند $A$ يساوي $(-e)$.

\textbf{II)} نعتبر الدالة العددية $g$ للمتغير الحقيقي $x$ المعرفة على المجال $]-2;+\infty[$ كما يلي: $g(x)=(-x-1)e^{-x}+1$، و $(C_g)$ تمثيلها البياني في نفس المعلم السابق.

\QQ{1} بيّن أن $\lim\limits_{x\to+\infty}g(x)=1$ وفسّر النتيجة بيانيًا. (نذكّر أن $\lim\limits_{u\to-\infty}ue^u=0$)
\QQ{2} أدرس تغيرات الدالة $g$، ثم أنشئ جدول تغيراتها.
\QQ{3} بيّن أن المنحنى $(C_g)$ يقبل نقطة انعطاف $I$ يطلب تعيين إحداثييها.
\QQ{4} أكتب معادلة المماس للمنحنى $(C_g)$ عند النقطة $I$.
\QQ{5} أرسم $(C_g)$.

\textbf{III)} لتكن $k$ الدالة المعرفة على المجال $]-2;+\infty[$ كما يأتي: $k(x)=g(x^2)$.

باستعمال مشتقة دالة مركبة، عيّن اتجاه تغير الدالة $k$ ثم شكّل جدول تغيراتها.
\end{tamrin}

\section{تمارين البكالوريا -- شعبة تقني رياضي}
% ===== تمارين البكالوريا في الدالة الأسية -- تقني رياضي (2026 -> 2018) =====
% graph frame with different units on the axes: {xmin}{xmax}{ymin}{ymax}{x unit}{y unit}
\newenvironment{GRbt}[6]{\par\begin{center}\begin{LTR}\begin{tikzpicture}[x=#5cm,y=#6cm,font=\footnotesize]
 \draw[gray!25,very thin] (#1,#3) grid (#2,#4);
 \draw[-{Stealth[length=2mm]}] (#1,0)--(#2+0.3,0) node[below left]{$x$};
 \draw[-{Stealth[length=2mm]}] (0,#3)--(0,#4+0.3) node[below left]{$y$};
 \node[below left] at (0,0){$O$};
 \begin{scope}\clip (#1,#3) rectangle (#2,#4);}
{\end{scope}\end{tikzpicture}\end{LTR}\end{center}}

%==================== 2026 ====================
\begin{tamrin}{بكالوريا 2026 -- الموضوع الأول -- تقني رياضي}
الدالة $f$ المعرفة على $\R$ بـ: $f(x)=1+\left(1-x-x^2\right)e^{-x}$، و$(C_f)$ تمثيلها البياني في المستوي المنسوب إلى المعلم المتعامد والمتجانس $(O;\vec{i},\vec{j})$.
\QQ{1} أحسب $\displaystyle\lim_{x\to-\infty}f(x)$، وبيّن أن $\displaystyle\lim_{x\to+\infty}f(x)=1$ ثم فسّر النتيجة هندسيا.
\QQ{2}
\begin{sq}
\qa بيّن أنه من أجل كل $x$ من $\R$: $f'(x)=\left(x^2-x-2\right)e^{-x}$.
\qb استنتج اتجاه تغير الدالة $f$ ثم شكّل جدول تغيراتها.
\end{sq}
\QQ{3} أثبت أن المعادلة $f(x)=0$ تقبل حلا وحيدا $\alpha$ في المجال $]-1{,}7;-1{,}6[$.
\QQ{4}
\begin{sq}
\qa عيّن معادلة $(T)$ مماس المنحنى $(C_f)$ عند النقطة ذات الفاصلة $0$.
\qb أحسب $f(-2)$ ثم ارسم $(T)$ و$(C_f)$.
\qc عيّن بيانيا قيم العدد الحقيقي $m$ حتى يكون للمعادلة $f(x)=m$ حلان سالبان تماما.
\end{sq}
\end{tamrin}

%==================== 2025 ====================
\begin{tamrin}{بكالوريا 2025 -- الموضوع الأول -- تقني رياضي}
\textbf{I)} $g$ الدالة المعرفة على $\R$ بـ: $g(x)=1-(x-1)^2e^{-x}$.
\QQ{1}
\begin{sq}
\qa أحسب $\displaystyle\lim_{x\to-\infty}g(x)$ ثم بيّن أن $\displaystyle\lim_{x\to+\infty}g(x)=1$.
\qb بيّن أنه من أجل كل عدد حقيقي $x$: $g'(x)=(x-1)(x-3)e^{-x}$.
\qc أدرس اتجاه تغير الدالة $g$ ثم شكّل جدول تغيراتها.
\end{sq}
\QQ{2} أحسب $g(0)$ ثم استنتج حسب قيم $x$ إشارة $g(x)$.

\textbf{II)} $f$ الدالة المعرفة على $\R$ بـ: $f(x)=x+\left(x^2+1\right)e^{-x}$، و$(C_f)$ تمثيلها البياني في المستوي المنسوب إلى المعلم المتعامد والمتجانس $(O;\vec{i},\vec{j})$ (وحدة الأطوال $2\,\text{cm}$).
\QQ{1}
\begin{sq}
\qa أحسب $\displaystyle\lim_{x\to-\infty}f(x)$ وبيّن أن $\displaystyle\lim_{x\to+\infty}f(x)=+\infty$.
\qb بيّن أن المستقيم $(\Delta)$ ذا المعادلة $y=x$ مقارب مائل لـ $(C_f)$ عند $+\infty$.
\end{sq}
\QQ{2}
\begin{sq}
\qa بيّن أنه من أجل كل عدد حقيقي $x$: $f'(x)=g(x)$.
\qb استنتج اتجاه تغير الدالة $f$ ثم شكّل جدول تغيراتها.
\end{sq}
\QQ{3}
\begin{sq}
\qa بيّن أن المنحنى $(C_f)$ يقبل مماسا $(T)$ موازيا لـ $(\Delta)$، يطلب تعيين معادلة له.
\qb تحقق أن للمنحنى $(C_f)$ نقطتي انعطاف، يطلب تعيين إحداثييهما.
\end{sq}
\QQ{4} أحسب $f(-1)$ ثم ارسم كلا من $(\Delta)$، $(T)$ و$(C_f)$.
\end{tamrin}

\begin{molahada}
السؤال \textbf{I) 2} (حساب $g(0)$ واستنتاج إشارة $g(x)$) سقط من نص المجلة، وهو ضروري لدراسة اتجاه تغير $f$؛ أعدناه كما في الموضوع الأصلي.
\end{molahada}

%==================== 2024 ====================
\begin{tamrin}{بكالوريا 2024 -- الموضوع الأول -- تقني رياضي}
\textbf{I)} $(C_h)$ التمثيل البياني للدالة $h$ المعرفة على $\R$ بـ: $h(x)=e^{-x}$، و$(D)$ المستقيم ذو المعادلة $y=x+2$، و$\alpha$ فاصلة نقطة تقاطع $(D)$ و$(C_h)$، كما في الشكل المقابل.
\begin{GRbt}{-2}{2.5}{-1}{4}{0.9}{0.9}
 \draw[cDef,very thick,domain=-1.4:2.5,samples=100] plot(\x,{exp(-\x)});
 \draw[cThm,thick,domain=-2:2] plot(\x,{\x+2});
 \draw[dashed] (-0.443,0)--(-0.443,1.557);
 \node[below,font=\scriptsize] at (-0.443,0){$\alpha$};
 \node[cDef] at (-1.1,3.5){$(C_h)$}; \node[cThm] at (1.2,3.6){$(D)$};
 \foreach \t in {-1,1,2} \node[below,font=\scriptsize] at (\t,0){$\t$};
 \foreach \t in {1,2,3} \node[left,font=\scriptsize] at (0,\t){$\t$};
\end{GRbt}
بقراءة بيانية: حدّد حسب قيم العدد الحقيقي $x$ إشارة $g(x)$ حيث: $g(x)=e^{-x}-x-2$.

\textbf{II)} $f$ الدالة العددية المعرفة على $\R$ بـ: $f(x)=-x+(x+1)e^{x}$، و$(C_f)$ تمثيلها البياني في المستوي المنسوب إلى المعلم المتعامد والمتجانس $(O;\vec{i},\vec{j})$ (وحدة الطول $2\,\text{cm}$).
\QQ{1} أحسب $\displaystyle\lim_{x\to-\infty}f(x)$ و$\displaystyle\lim_{x\to+\infty}f(x)$.
\QQ{2}
\begin{sq}
\qa بيّن أنه من أجل كل عدد حقيقي $x$: $f'(x)=-g(x)e^{x}$.
\qb استنتج اتجاه تغير الدالة $f$، ثم شكّل جدول تغيراتها.
\end{sq}
\QQ{3}
\begin{sq}
\qa بيّن أن المستقيم $(\Delta)$ ذا المعادلة $y=-x$ مقارب مائل للمنحنى $(C_f)$ عند $-\infty$.
\qb أدرس الوضع النسبي للمنحنى $(C_f)$ والمستقيم $(\Delta)$.
\end{sq}
\QQ{4} بيّن أن $(C_f)$ يقبل مماسا $(T)$ موازيا لـ $(\Delta)$، يطلب تعيين معادلة له.
\QQ{5}
\begin{sq}
\qa أرسم $(T)$، $(\Delta)$ و$(C_f)$ (نأخذ: $\alpha\approx-0{,}44$ و$f(\alpha)\approx0{,}8$).
\qb عيّن بيانيا قيم الوسيط الحقيقي $m$ التي من أجلها تقبل المعادلة $f(x)=-x-e^{m}$ حلين مختلفين.
\end{sq}
\end{tamrin}

\begin{molahada}
في نص المجلة: «$\alpha\approx0{,}5$»، وهذا خطأ لأن $\alpha$ سالب (الشكل، و$g(-0{,}5)\approx0{,}15>0$، $g(-0{,}4)\approx-0{,}11<0$)؛ الصواب $\alpha\approx-0{,}44$.
\end{molahada}

%==================== 2023 ====================
\begin{tamrin}{بكالوريا 2023 -- الموضوع الثاني -- تقني رياضي}
\textbf{I)} الجدول التالي يمثل تغيرات الدالة $g$ المعرفة على $\R$ بـ: $g(x)=-1+(2x-1)e^{x}$.
\begin{TAB}
\tkzTabInit[lgt=2,espcl=3.5]{$x$ /0.8, $g(x)$ /1.8}{$-\infty$, $-\frac12$, $+\infty$}
\tkzTabVar{+/$-1$, -/$g\left(-\frac12\right)$, +/$+\infty$}
\end{TAB}
\QQ{1} أثبت أن المعادلة $g(x)=0$ تقبل حلا وحيدا $\alpha$ حيث $0{,}7<\alpha<0{,}8$.
\QQ{2} استنتج حسب قيم $x$ إشارة $g(x)$ على $\R$.

\textbf{II)} $f$ الدالة المعرفة على $\R$ بـ: $f(x)=-x+4+(2x-3)e^{x}$، و$(C_f)$ تمثيلها البياني في المستوي المنسوب إلى المعلم المتعامد والمتجانس $(O;\vec{i},\vec{j})$.
\QQ{1}
\begin{sq}
\qa أحسب $\displaystyle\lim_{x\to-\infty}f(x)$ ثم بيّن أن $\displaystyle\lim_{x\to+\infty}f(x)=+\infty$.
\qb بيّن أن المستقيم $(\Delta)$ ذا المعادلة $y=-x+4$ مقارب مائل لـ $(C_f)$ عند $-\infty$.
\qc أدرس وضعية $(C_f)$ بالنسبة إلى $(\Delta)$.
\end{sq}
\QQ{2}
\begin{sq}
\qa بيّن أنه من أجل كل عدد حقيقي $x$: $f'(x)=g(x)$.
\qb استنتج أن الدالة $f$ متناقصة تماما على المجال $]{-}\infty;\alpha]$ ومتزايدة تماما على المجال $[\alpha;+\infty[$ ثم شكّل جدول تغيراتها.
\end{sq}
\QQ{3}
\begin{sq}
\qa أثبت أن $(C_f)$ يقبل مماسا $(T)$ موازيا لـ $(\Delta)$، يطلب تعيين معادلة له.
\qb أرسم $(\Delta)$، $(T)$ و$(C_f)$ (نأخذ: $f(2)\approx9{,}4$ و$f(\alpha)\approx0{,}1$).
\qc عيّن بيانيا قيم الوسيط الحقيقي $m$ التي من أجلها تقبل المعادلة $f(x)=-x+m$ حلين بالضبط.
\end{sq}
\end{tamrin}

%==================== 2022 ====================
\begin{tamrin}{بكالوريا 2022 -- الموضوع الثاني -- تقني رياضي}
\textbf{I)} المستوي منسوب إلى المعلم المتعامد والمتجانس $(O;\vec{i},\vec{j})$. $(\Gamma)$ التمثيل البياني للدالة $x\mapsto e^{-x}$ و$(C_g)$ التمثيل البياني للدالة $g$ المعرفة على $\R$ بـ: $g(x)=\dfrac{x^2+2x-1}{\left(x^2+1\right)^2}$، و$\alpha$ فاصلة نقطة تقاطع $(C_g)$ و$(\Gamma)$ (كما هو مبيّن في الشكل المقابل).
\begin{GRbt}{-2}{3.5}{-1.5}{3}{1}{1}
 \draw[cThm,very thick,domain=-1.1:3.5,samples=100] plot(\x,{exp(-\x)});
 \draw[cDef,very thick,domain=-2:3.5,samples=200] plot(\x,{((\x)*(\x)+2*\x-1)/(((\x)*(\x)+1)*((\x)*(\x)+1))});
 \draw[dashed] (0.787,0)--(0.787,0.455);
 \node[below,font=\scriptsize] at (0.787,0){$\alpha$};
 \node[cThm] at (-0.6,2.6){$(\Gamma)$}; \node[cDef] at (-1.4,-0.5){$(C_g)$};
 \foreach \t in {-1,1,2,3} \node[below,font=\scriptsize] at (\t,0){$\t$};
 \foreach \t in {-1,1,2} \node[left,font=\scriptsize] at (0,\t){$\t$};
\end{GRbt}
\QQ{1} بقراءة بيانية، حدّد حسب قيم العدد الحقيقي $x$ إشارة $g(x)-e^{-x}$.
\QQ{2} تحقق حسابيا أن $0{,}7<\alpha<0{,}8$.

\textbf{II)} $f$ الدالة العددية المعرفة على $\R$ بـ: $f(x)=e^{-x}-\dfrac{x+1}{x^2+1}$، و$(C_f)$ تمثيلها البياني في معلم متعامد ومتجانس.
\QQ{1} أحسب $\displaystyle\lim_{x\to-\infty}f(x)$ ثم أحسب $\displaystyle\lim_{x\to+\infty}f(x)$ وفسّر النتيجة بيانيا.
\QQ{2}
\begin{sq}
\qa بيّن أنه من أجل كل عدد حقيقي $x$: $f'(x)=g(x)-e^{-x}$.
\qb استنتج اتجاه تغير الدالة $f$، ثم شكّل جدول تغيراتها.
\end{sq}
\QQ{3}
\begin{sq}
\qa أحسب $\displaystyle\lim_{x\to-\infty}\left[f(x)-e^{-x}\right]$ وفسّر النتيجة بيانيا.
\qb أدرس الوضعية النسبية للمنحنيين $(C_f)$ و$(\Gamma)$.
\end{sq}
\QQ{4}
\begin{sq}
\qa أكتب معادلة لـ $(T)$ مماس المنحنى $(C_f)$ في النقطة ذات الفاصلة $0$.
\qb أنشئ $(T)$، $(\Gamma)$ و$(C_f)$.
\qc ناقش بيانيا وحسب قيم الوسيط الحقيقي $m$ عدد وإشارة حلول المعادلة $f(x)-m=0$.
\end{sq}
\end{tamrin}

%==================== 2021 ====================
\begin{tamrin}{بكالوريا 2021 -- الموضوع الأول -- تقني رياضي}
\textbf{I)} الدالة العددية $g$ معرفة على المجال $[0;+\infty[$ بـ: $g(x)=x^2-5+e^{x-1}$.
\QQ{1} بيّن أن الدالة $g$ متزايدة تماما على المجال $[0;+\infty[$.
\QQ{2}
\begin{sq}
\qa بيّن أن المعادلة $g(x)=0$ تقبل حلا وحيدا $\alpha$ حيث: $1{,}71<\alpha<1{,}72$.
\qb استنتج حسب قيم العدد الحقيقي الموجب $x$ إشارة $g(x)$.
\end{sq}

\textbf{II)} الدالة العددية $f$ معرفة على المجال $[0;+\infty[$ بـ: $f(x)=x+1+\left(-x^2-2x+3\right)e^{1-x}$، و$(C_f)$ تمثيلها البياني في المستوي المنسوب إلى المعلم المتعامد والمتجانس $(O;\vec{i},\vec{j})$.
\QQ{1}
\begin{sq}
\qa بيّن أنه من أجل كل عدد حقيقي $x$ من المجال $[0;+\infty[$: $f'(x)=g(x)e^{1-x}$.
\qb استنتج أن الدالة $f$ متزايدة تماما على المجال $[\alpha;+\infty[$ ومتناقصة تماما على المجال $[0;\alpha]$.
\qc بيّن أن $\displaystyle\lim_{x\to+\infty}f(x)=+\infty$، ثم شكّل جدول تغيرات الدالة $f$.
\end{sq}
\QQ{2} بيّن أن المستقيم $(\Delta)$ ذا المعادلة $y=x+1$ مقارب مائل لـ $(C_f)$، ثم أدرس وضعية $(C_f)$ بالنسبة لـ $(\Delta)$.
\QQ{3} بيّن أن $(C_f)$ يقبل مماسا $(T)$ موازيا لـ $(\Delta)$ في نقطة $A$ يطلب تعيين فاصلتها (لا يطلب كتابة معادلة $(T)$).
\QQ{4}
\begin{sq}
\qa بيّن أن $(C_f)$ يقبل نقطة انعطاف وحيدة فاصلتها $1+\sqrt6$.
\qb أرسم $(\Delta)$، $(T)$ و$(C_f)$ (نأخذ: $f(\alpha)\approx1{,}1$، $f\left(\sqrt5\right)\approx1{,}4$ و$f\left(1+\sqrt6\right)\approx3{,}1$).
\end{sq}
\QQ{5} الدالة العددية $h$ معرفة على المجال $]{-}\infty;0]$ بـ: $h(x)=-x+1+\left(-x^2+2x+3\right)e^{1+x}$، و$(C_h)$ تمثيلها البياني في المعلم السابق.
\begin{sq}
\qa تحقق أنه من أجل كل عدد حقيقي $x$ من المجال $]{-}\infty;0]$: $h(x)=f(-x)$.
\qb إشرح كيفية رسم $(C_h)$ انطلاقا من $(C_f)$ ثم ارسمه.
\end{sq}
\end{tamrin}

%==================== 2020 ====================
\begin{tamrin}{بكالوريا 2020 -- الموضوع الثاني -- تقني رياضي}
نعتبر الدالة العددية $f$ المعرفة على المجال $[-1;+\infty[$ كما يلي: $f(x)=x-1+\dfrac14\left(2e^{-x}-1\right)^2$، و$(C_f)$ تمثيلها البياني في المستوي المنسوب إلى المعلم المتعامد والمتجانس $(O;\vec{i},\vec{j})$ (وحدة الطول $2\,\text{cm}$).
\QQ{1}
\begin{sq}
\qa بيّن أنه من أجل كل $x$ من المجال $[-1;+\infty[$: $f'(x)=\left(1-e^{-x}\right)\left(2e^{-x}+1\right)$.
\qb أدرس إشارة $f'(x)$، واستنتج اتجاه تغير الدالة $f$.
\qc أحسب $\displaystyle\lim_{x\to+\infty}f(x)$، ثم شكّل جدول تغيراتها.
\end{sq}
\QQ{2}
\begin{sq}
\qa بيّن أن المستقيم $(\Delta)$ ذا المعادلة $y=x-\dfrac34$ مقارب مائل للمنحنى $(C_f)$.
\qb أدرس وضعية المنحنى $(C_f)$ بالنسبة للمستقيم $(\Delta)$.
\end{sq}
\QQ{3} بيّن أن المنحنى $(C_f)$ يقبل مماسا $(T)$ يوازي $(\Delta)$، يطلب كتابة معادلة له.
\QQ{4} بيّن أن المنحنى $(C_f)$ يقبل نقطة انعطاف يطلب تعيينها.
\QQ{5} أرسم $(\Delta)$، $(T)$ و$(C_f)$.
\QQ{6} ليكن $m$ وسيطا حقيقيا. عيّن مجموعة قيم $m$ التي من أجلها تقبل المعادلة $f(x)=x+m$ حلين مختلفين.
\end{tamrin}

%==================== 2019 ====================
\Needspace{18\baselineskip}
\begin{tamrin}{بكالوريا 2019 -- الموضوع الأول -- تقني رياضي}
\textbf{I)} $g$ الدالة العددية المعرفة على $\R$ كما يلي: $g(x)=(x+3)e^{x}-1$، و$(C_g)$ تمثيلها البياني كما هو مبيّن في الشكل.
\begin{GRbt}{-3}{1.5}{-1.5}{5}{0.9}{0.6}
 \draw[cDef,very thick,domain=-3:0.6,samples=150] plot(\x,{(\x+3)*exp(\x)-1});
 \node[cDef] at (0.9,3.8){$(C_g)$};
 \foreach \t in {-2,-1,1} \node[below,font=\scriptsize] at (\t,0){$\t$};
 \foreach \t in {-1,1,2,3,4} \node[left,font=\scriptsize] at (0,\t){$\t$};
\end{GRbt}
بقراءة بيانية:
\QQ{1} حدّد إشارة $g(-1)$ و$g\left(-\frac12\right)$.
\QQ{2} استنتج وجود عدد حقيقي $\alpha$ وحيد من المجال $\left]-1;-\frac12\right[$ بحيث: $g(\alpha)=0$، ثم تحقق أن: $-0{,}8<\alpha<-0{,}7$.
\QQ{3} استنتج إشارة $g(x)$ على $\R$.

\textbf{II)} نعتبر الدالة العددية $f$ المعرفة على $\R$ كما يلي: $f(x)=(x+2)\left(e^{x}-1\right)$، و$(C_f)$ تمثيلها البياني في المستوي المنسوب إلى المعلم المتعامد والمتجانس $(O;\vec{i},\vec{j})$.
\QQ{1} أحسب $\displaystyle\lim_{x\to-\infty}f(x)$ و$\displaystyle\lim_{x\to+\infty}f(x)$.
\QQ{2} بيّن أنه من أجل كل عدد حقيقي $x$: $f'(x)=g(x)$، ثم شكّل جدول تغيراتها.
\QQ{3}
\begin{sq}
\qa أحسب $\displaystyle\lim_{x\to-\infty}\left[f(x)+x\right]$ ثم استنتج أن $(C_f)$ يقبل مستقيما مقاربا مائلا $(\Delta)$ يطلب تعيين معادلة له.
\qb أدرس وضعية المنحنى $(C_f)$ بالنسبة للمستقيم $(\Delta)$.
\qc أكتب معادلة لـ $(T)$ مماس $(C_f)$ الموازي للمستقيم $(\Delta)$.
\end{sq}
\QQ{4} أرسم المستقيم $(\Delta)$ والمنحنى $(C_f)$ على المجال $]{-}\infty;1]$ (يعطى $f(\alpha)\approx-0{,}7$).
\QQ{5} $h$ الدالة المعرفة على $\R$ كما يلي: $h(x)=|x|\left(e^{|x|-2}-1\right)+1$، و$(C_h)$ تمثيلها البياني في المعلم السابق.
\begin{sq}
\qa بيّن أن الدالة $h$ زوجية.
\qb تحقق أنه من أجل كل $x$ من المجال $[0;+\infty[$ فإن: $h(x)=f(x-2)+1$.
\qc إشرح كيف يمكن رسم $(C_h)$ انطلاقا من $(C_f)$ ثم ارسم $(C_h)$ على المجال $[-3;3]$.
\end{sq}
\end{tamrin}

%==================== 2018 ====================
\begin{tamrin}{بكالوريا 2018 -- الموضوع الأول -- تقني رياضي}
$f$ الدالة العددية المعرفة على المجال $]{-}\infty;1[$ بـ: $f(x)=\dfrac{x}{x-1}e^{-x}$، و$(C_f)$ تمثيلها البياني في المستوي المنسوب إلى المعلم المتعامد والمتجانس $(O;\vec{i},\vec{j})$.
\QQ{1} أحسب $\displaystyle\lim_{\substack{x\to1\\x<1}}f(x)$ ثم فسّر النتيجة بيانيا، وأحسب $\displaystyle\lim_{x\to-\infty}f(x)$.
\QQ{2} بيّن أنه من أجل كل $x$ من $]{-}\infty;1[$: $f'(x)=\dfrac{\left(-x^2+x-1\right)e^{-x}}{(x-1)^2}$، وأدرس اتجاه تغير الدالة $f$ ثم شكّل جدول تغيراتها.
\QQ{3}
\begin{sq}
\qa أكتب معادلة المماس $(T)$ للمنحنى $(C_f)$ عند النقطة ذات الفاصلة $0$.
\qb $h$ الدالة المعرفة على المجال $]{-}\infty;1[$ بـ: $h(x)=e^{-x}+x-1$. أدرس اتجاه تغير الدالة $h$، ثم استنتج أنه من أجل كل $x$ من $]{-}\infty;1[$: $h(x)\geqslant0$.
\end{sq}
\QQ{4} بيّن أنه من أجل كل $x$ من $]{-}\infty;1[$: $f(x)+x=\dfrac{xh(x)}{x-1}$، ثم استنتج الوضع النسبي للمنحنى $(C_f)$ والمماس $(T)$. فسّر النتيجة بيانيا.
\QQ{5} أكتب معادلة المستقيم $(\Delta)$ الذي يشمل مبدأ المعلم $O$ والنقطة $A\left(-2;\frac23e^2\right)$، ثم أرسم المستقيمين $(\Delta)$ و$(T)$ والمنحنى $(C_f)$ على المجال $[-2;1[$.
\QQ{6} $m$ وسيط حقيقي، ناقش بيانيا وحسب قيم الوسيط الحقيقي $m$ عدد حلول المعادلة: $f(x)=mx$، حيث $x\in[-2;1[$.
\end{tamrin}

%====================================================== 2017 استثنائية
\begin{tamrin}{بكالوريا 2017 -- الموضوع الثاني -- تقني رياضي -- الدورة الاستثنائية}
\textbf{I)} $g$ الدالة العددية المعرفة على $\R$ كما يلي: $g(x)=1-2xe^{-x}$.
\par ادرس اتجاه تغير الدالة $g$، ثم استنتج إشارة $g(x)$.

\textbf{II)} نعتبر الدالة العددية $f$ المعرفة على $\R$ كما يلي: $f(x)=(x+1)\left(1+2e^{-x}\right)$،
و $(C_f)$ تمثيلها البياني في المستوي المنسوب إلى المعلم المتعامد المتجانس $(O;\vec i,\vec j)$.
\QQ{1}
\begin{sq}
\qa احسب $\displaystyle\lim_{x\to-\infty}f(x)$ و $\displaystyle\lim_{x\to+\infty}f(x)$.
\qb ادرس اتجاه تغير الدالة $f$، ثم شكّل جدول تغيراتها.
\end{sq}
\QQ{2}
\begin{sq}
\qa بيّن أن $\displaystyle\lim_{x\to+\infty}\big[f(x)-x\big]=1$، ثم استنتج معادلة لـ $(\Delta)$، المستقيم المقارب المائل للمنحنى $(C_f)$.
\qb ادرس وضعية المنحنى $(C_f)$ بالنسبة للمستقيم $(\Delta)$.
\end{sq}
\QQ{3} أثبت أن المنحنى $(C_f)$ يقبل مماسا وحيدا $(T)$ يوازي $(\Delta)$ يطلب تعيين معادلة له.
\QQ{4} باستعمال المنحنى $(C_f)$، عيّن قيم الوسيط الحقيقي $m$ حتى يكون للمعادلة $f(x)=x+m$ حلان مختلفان.
\end{tamrin}

%====================================================== 2015
\begin{tamrin}{بكالوريا 2015 -- الموضوع الثاني -- تقني رياضي}
\textbf{I)} $g$ الدالة العددية المعرفة على $\R$ كما يلي: $g(x)=(x+2)e^{x}-2$.
\QQ{1} احسب $\displaystyle\lim_{x\to-\infty}g(x)$ و $\displaystyle\lim_{x\to+\infty}g(x)$.
\QQ{2} ادرس اتجاه تغير الدالة $g$ ثم شكّل جدول تغيراتها.
\QQ{3} احسب $g(0)$، ثم استنتج إشارة $g(x)$.

\textbf{II)} نعتبر الدالة $f$ المعرفة على $\R$ بـ: $f(x)=2x+3-(x+1)e^{x}$،
و $(C_f)$ تمثيلها البياني في المستوي المنسوب إلى المعلم المتعامد المتجانس $(O;\vec i,\vec j)$.
\QQ{1} بيّن أن $\displaystyle\lim_{x\to+\infty}f(x)=-\infty$، ثم احسب $\displaystyle\lim_{x\to-\infty}f(x)$.
\QQ{2}
\begin{sq}
\qa بيّن أنه من أجل كل عدد حقيقي $x$: $f'(x)=-g(x)$.
\qb استنتج إشارة $f'(x)$، ثم شكّل جدول تغيرات الدالة $f$.
\qc بيّن أن المستقيم $(\Delta)$ ذا المعادلة $y=2x+3$ مستقيم مقارب مائل للمنحنى $(C_f)$ عند $-\infty$، ثم ادرس وضعية $(C_f)$ بالنسبة للمستقيم $(\Delta)$.
\end{sq}
\QQ{3}
\begin{sq}
\qa بيّن أن المعادلة $f(x)=0$ تقبل حلين $\alpha$ و $\beta$ حيث: $0{,}92<\alpha<0{,}93$ و $-1{,}56<\beta<-1{,}55$.
\qb أنشئ المستقيم $(\Delta)$ والمنحنى $(C_f)$ على المجال $\left]{-}\infty;\frac32\right]$.
\end{sq}
\end{tamrin}

%====================================================== 2014
\begin{tamrin}{بكالوريا 2014 -- الموضوع الثاني -- تقني رياضي}
$f$ هي الدالة المعرفة على $\R$ بـ: $f(x)=(x-1)e^{x}$،
و $(C_f)$ تمثيلها البياني في المستوي المنسوب إلى المعلم المتعامد المتجانس $(O;\vec i,\vec j)$.
\QQ{1} عيّن نهاية $f$ عند كل من $+\infty$ و $-\infty$.
\QQ{2} ادرس اتجاه تغير الدالة $f$ على $\R$، ثم شكّل جدول تغيراتها.
\QQ{3}
\begin{sq}
\qa بيّن أن المعادلة $f(x)=1$ تقبل حلا وحيدا $\alpha$ على $\R$، ثم تحقق أن: $1{,}27<\alpha<1{,}28$.
\qb اكتب معادلة لـ $(T)$ مماس المنحنى $(C_f)$ عند النقطة ذات الفاصلة $1$، وحدد وضعية $(C_f)$ بالنسبة لـ $(T)$.
\qc ارسم $(T)$ و $(C_f)$.
\end{sq}
\QQ{4} عيّن قيم العدد الحقيقي $m$ التي من أجلها تقبل المعادلة: $(x-1)e^{x}-(m-1)e^{m}=-1$ حلا واحدا في $\R$.
\QQ{5} $h$ هي الدالة المعرفة على $\R$ بـ: $h(x)=\left(|x|+1\right)e^{-|x|}$، و $(C_h)$ تمثيلها البياني.
\begin{sq}
\qa بيّن أن الدالة $h$ زوجية.
\qb ارسم $(C_h)$ مستعينا بالمنحنى $(C_f)$.
\end{sq}
\QQ{6} $g$ دالة معرفة على $\R$ بـ: $g(x)=(ax+b)e^{x}$، حيث $a$ و $b$ عددان حقيقيان.
\par عيّن $a$ و $b$ حتى يكون من أجل كل $x$ من $\R$: $g'(x)=f(x)$.
\end{tamrin}

%====================================================== 2012
\begin{tamrin}{بكالوريا 2012 -- الموضوع الأول -- تقني رياضي}
\textbf{I)} $g$ هي الدالة العددية المعرفة على $\R$ بـ: $g(x)=-4+(4-2x)e^{x}$.
\QQ{1} ادرس تغيرات الدالة $g$، ثم شكّل جدول تغيراتها.
\QQ{2} بيّن أن المعادلة $g(x)=0$ تقبل حلين أحدهما معدوم والآخر $\alpha$ حيث: $1{,}59<\alpha<1{,}60$.
\QQ{3} استنتج إشارة $g(x)$.

\textbf{II)} $f$ هي الدالة المعرفة على $\R$ كما يلي: $f(x)=\dfrac{2x-2}{e^{x}-2x}$،
و $(C_f)$ تمثيلها البياني في المستوي المنسوب إلى المعلم المتعامد المتجانس $(O;\vec i,\vec j)$ (وحدة الطول $2\,\mathrm{cm}$).
\QQ{1} بيّن أن $(C_f)$ يقبل عند $-\infty$ و $+\infty$ مستقيمين مقاربين معادلتاهما على الترتيب $y=-1$ و $y=0$.
\QQ{2}
\begin{sq}
\qa برهن أنه من أجل كل عدد حقيقي $x$: $f'(x)=\dfrac{g(x)}{\left(e^{x}-2x\right)^2}$.
\qb استنتج إشارة $f'(x)$، ثم شكّل جدول تغيرات الدالة $f$.
\qc احسب $f(1)$، ثم استنتج، حسب قيم $x$، إشارة $f(x)$.
\end{sq}
\QQ{3}
\begin{sq}
\qa بيّن أن $f(\alpha)=-1+\dfrac{1}{\alpha-1}$، حيث $\alpha$ هو العدد المعرف في السؤال 2 من الجزء I.
\qb استنتج حصرا للعدد $f(\alpha)$ (تدوّر النتائج إلى $10^{-2}$).
\qc ارسم $(C_f)$.
\end{sq}
\QQ{4} ناقش بيانيا، حسب قيم الوسيط الحقيقي $m$، عدد وإشارة حلول المعادلة: $2x-2=\left(e^{x}-2x\right)(m+1)$.
\QQ{5} $h$ هي الدالة المعرفة على $\R$ كما يلي: $h(x)=\big[f(x)\big]^2$.
\begin{sq}
\qa احسب $h'(x)$ بدلالة كل من $f(x)$ و $f'(x)$، ثم استنتج إشارة $h'(x)$.
\qb شكّل جدول تغيرات الدالة $h$.
\end{sq}
\end{tamrin}

%====================================================== 2011
\begin{tamrin}{بكالوريا 2011 -- الموضوع الثاني -- تقني رياضي}
$f$ الدالة المعرفة على $\R$ كما يلي: $f(x)=3-\dfrac{4}{e^{x}+1}$،
و $(C_f)$ تمثيلها البياني في المستوي المنسوب إلى المعلم المتعامد المتجانس $(O;\vec i,\vec j)$.
\QQ{1} ادرس تغيرات الدالة $f$.
\QQ{2} عيّن المستقيمات المقاربة للمنحنى $(C_f)$.
\QQ{3} بيّن أن للمنحنى $(C_f)$ نقطة انعطاف $\omega$ يطلب تعيينها، ثم اكتب معادلة لمماس $(C_f)$ عندها.
\QQ{4} لتكن $g$ الدالة العددية المعرفة على $\R$ كما يلي: $g(x)=f(x)-x$.
\begin{sq}
\qa ادرس تغيرات الدالة $g$.
\qb بيّن أن المعادلة $g(x)=0$ تقبل حلا وحيدا $\alpha$ حيث $2{,}7<\alpha<2{,}8$.
\end{sq}
\QQ{5}
\begin{sq}
\qa حل في $\R$ المعادلة: $f(x)=0$.
\qb ارسم المماس والمستقيم $(\Delta)$ الذي معادلته $y=x$ والمنحنى $(C_f)$.
\end{sq}
\end{tamrin}

%====================================================== 2010
\begin{tamrin}{بكالوريا 2010 -- الموضوع الثاني -- تقني رياضي}
$f$ الدالة العددية المعرفة على $\R^*$ كما يلي: $f(x)=\dfrac{3xe^{x}-3x-4}{3\left(e^{x}-1\right)}$،
و $(C_f)$ تمثيلها البياني في المستوي المنسوب إلى المعلم المتعامد المتجانس $(O;\vec i,\vec j)$.
\QQ{1} عيّن العددين الحقيقيين $a$ و $b$ بحيث من أجل كل $x$ من $\R^*$: $f(x)=ax+\dfrac{b}{3\left(e^{x}-1\right)}$.
\QQ{2} احسب نهايات الدالة $f$ عند أطراف مجموعة تعريفها.
\QQ{3} بيّن أن $f$ متزايدة تماما على كل مجال من مجالي تعريفها، ثم شكّل جدول تغيراتها.
\QQ{4}
\begin{sq}
\qa $(D)$ و $(D')$ المستقيمان اللذان معادلتاهما على الترتيب: $y=x+\dfrac43$ و $y=x$.
\par بيّن أن $(D)$ و $(D')$ مقاربان للمنحنى $(C_f)$، ثم حدد وضعيته بالنسبة لكل منهما.
\qb بيّن أن المعادلة $f(x)=0$ تقبل حلين $x_0$ و $x_1$ حيث: $0{,}9<x_0<0{,}91$ و $-1{,}66<x_1<-1{,}65$.
\qc احسب من أجل كل عدد حقيقي $x$ غير معدوم $f(x)+f(-x)$. فسر النتيجة هندسيا.
\qd ارسم $(D)$ و $(D')$ و $(C_f)$.
\qe $m$ عدد حقيقي، $(D_m)$ المستقيم المعرف بالمعادلة $y=x+m$.
\par ناقش بيانيا حسب قيم $m$ عدد حلول المعادلة: $f(x)=x+m$.
\end{sq}
\QQ{5} نعتبر الدالة $g$ المعرفة على المجال $]0;+\infty[$ كما يأتي: $g(x)=\big[f(x)\big]^2$.
\par ادرس تغيرات الدالة $g$ دون حساب $g(x)$ بدلالة $x$.
\end{tamrin}

%====================================================== 2009
\begin{tamrin}{بكالوريا 2009 -- الموضوع الأول -- تقني رياضي}
نعتبر الدالة $f$ المعرفة على $\R$ كما يلي: $f(x)=x+\dfrac{2}{e^{x}+1}$،
و $(C_f)$ تمثيلها البياني في المستوي المنسوب إلى المعلم المتعامد المتجانس $(O;\vec i,\vec j)$.
\QQ{1} احسب من أجل كل عدد حقيقي $x$: $f(x)+f(-x)$، ثم استنتج أن النقطة $\omega(0;1)$ هي مركز تناظر للمنحنى $(C_f)$.
\QQ{2} ادرس تغيرات الدالة $f$ على المجال $[0;+\infty[$، ثم استنتج جدول تغيراتها على $\R$.
\QQ{3}
\begin{sq}
\qa بيّن أن المستقيم ذا المعادلة $y=x$ مقارب مائل للمنحنى $(C_f)$ عند $+\infty$.
\qb احسب $\displaystyle\lim_{x\to-\infty}\big[f(x)-(x+2)\big]$، استنتج المستقيم المقارب للمنحنى $(C_f)$ عند $-\infty$.
\end{sq}
\QQ{4} بيّن أن للمعادلة $f(x)=0$ حلا وحيدا $\alpha$ بحيث: $-1{,}7<\alpha<-1{,}6$.
\QQ{5} ارسم $(C_f)$.
\end{tamrin}

\section{تمارين البكالوريا -- شعبة رياضيات}
% ===================== بكالوريا 2026 -- الموضوع الثاني -- رياضيات =====================
\Needspace{16\baselineskip}
\begin{tamrin}{بكالوريا 2026 -- الموضوع الثاني -- رياضيات}
\textbf{I)} الدالة $g$ المعرفة على $\R$ بـ: $g(x)=(x^2+2x)e^{x}+2xe^{2x}$.

التمثيل البياني $(C_g)$ للدالة $g$ يشمل النقطة $A(\alpha;0)$ كما هو موضح في الشكل المقابل.
\begin{GR}{-4.5}{1.5}{-1.5}{3.5}{1}
 \Ticks{-4,-3,-2,-1,1}{-1,1,2,3}
 \draw[cDef,very thick,domain=-4.5:0.42,samples=200,smooth] plot (\x,{(\x*\x+2*\x)*exp(\x)+2*\x*exp(2*\x)});
 \fill[cThm] (-2.2177,0) circle (1.8pt) node[above]{$A$};
 \node[cDef] at (-3.3,0.45){$(C_g)$};
\end{GR}
\QQ{1} تحقق أن: $-2{,}3<\alpha<-2{,}2$.
\QQ{2} بقراءة بيانية، حدد حسب قيم العدد الحقيقي $x$ إشارة $g(x)$.

\textbf{II)} الدالة $f$ المعرفة على $\R$ بـ: $f(x)=\dfrac{x^2e^{x}}{e^{x}+1}$، و $(C_f)$ تمثيلها البياني في المستوي المنسوب إلى المعلم المتعامد والمتجانس $(O;\vec i,\vec j)$.
\QQ{1} احسب $\lim\limits_{x\to-\infty}f(x)$ وبين أن: $\lim\limits_{x\to+\infty}f(x)=+\infty$.
\QQ{2} بين أنه من أجل كل $x$ من $\R$: $f'(x)=\dfrac{g(x)}{(e^{x}+1)^2}$ ثم شكل جدول تغيرات الدالة $f$.
\QQ{3}
\begin{sq}
\qa بين أن المنحنى $(\Gamma)$ ذا المعادلة $y=x^2$ مقارب للمنحنى $(C_f)$ عند $+\infty$.
\qb ادرس الوضع النسبي للمنحنيين $(C_f)$ و $(\Gamma)$ ثم ارسم $(\Gamma)$ و $(C_f)$. (نأخذ $f(\alpha)\simeq 0{,}5$).
\end{sq}
\end{tamrin}

% ===================== بكالوريا 2025 -- الموضوع الأول -- رياضيات =====================
\Needspace{16\baselineskip}
\begin{tamrin}{بكالوريا 2025 -- الموضوع الأول -- رياضيات}
\textbf{I)} \QQ{1} $g$ الدالة المعرفة على $\R$ بـ: $g(x)=1+xe^{x}$.

أدرس اتجاه تغير الدالة $g$ ثم استنتج أنه من أجل كل عدد حقيقي $x$: $1+xe^{x}>0$.
\QQ{2} $h$ الدالة المعرفة على $\R$ بـ: $h(x)=(x^2+1)e^{2x}+(x-2)e^{x}+1$.
تمثيلها البياني $(C_h)$ يقطع حامل محور الفواصل في نقطتين فاصلتاهما $\alpha$ و $0$ كما في الشكل.
\begin{GR}{-4}{1.2}{-0.6}{2.5}{1.2}
 \Ticks{-3,-2,-1,1}{1,2}
 \draw[cDef,very thick,domain=-4:0.42,samples=200,smooth] plot (\x,{(\x*\x+1)*exp(2*\x)+(\x-2)*exp(\x)+1});
 \fill[cThm] (-0.6403,0) circle (1.6pt) node[above left]{$\alpha$};
 \node[cDef] at (-2.8,1.25){$(C_h)$};
\end{GR}
\begin{sq}
\qa بقراءة بيانية، حدد إشارة $h(x)$ على $\R$.
\qb تحقق أن: $-0{,}7<\alpha<-0{,}6$.
\end{sq}

\textbf{II)} $f$ الدالة المعرفة على $\R$ بـ: $f(x)=2x+\dfrac{2-2e^{x}}{1+xe^{x}}$، و $(C_f)$ تمثيلها البياني في المستوي المنسوب إلى المعلم المتعامد والمتجانس $(O;\vec i,\vec j)$ (وحدة الطول $2cm$).
\QQ{1}
\begin{sq}
\qa احسب $\lim\limits_{x\to-\infty}f(x)$ وبين أن $\lim\limits_{x\to+\infty}f(x)=+\infty$.
\qb بين أن المستقيم $(\Delta)$ ذا المعادلة $y=2x+2$ مقارب مائل لـ $(C_f)$ عند $-\infty$، وأن المستقيم $(\Delta')$ ذا المعادلة $y=2x$ مقارب مائل لـ $(C_f)$ عند $+\infty$.
\qc ادرس وضعية $(C_f)$ بالنسبة إلى كل من $(\Delta)$ و $(\Delta')$.
\end{sq}
\QQ{2}
\begin{sq}
\qa بين أنه من أجل كل عدد حقيقي $x$: $f'(x)=\dfrac{2h(x)}{\left(1+xe^{x}\right)^2}$.
\qb استنتج اتجاه تغير الدالة $f$ ثم شكل جدول تغيراتها.
\end{sq}
\QQ{3} أرسم كلا من $(\Delta)$، $(\Delta')$ و $(C_f)$. (نأخذ $f(\alpha)\simeq0{,}15$).
\QQ{4} $k$ الدالة المعرفة على $\R$ بـ: $k(x)=f\left(x^2\right)$.
دون حساب عبارة $k(x)$، حدد اتجاه تغير الدالة $k$ ثم شكل جدول تغيراتها.
\end{tamrin}

% ===================== بكالوريا 2024 -- الموضوع الأول -- رياضيات =====================
\Needspace{16\baselineskip}
\begin{tamrin}{بكالوريا 2024 -- الموضوع الأول -- رياضيات}
\textbf{I)} يمثل الجدول المقابل تغيرات الدالة $g$ المعرفة على $\R$ بـ: $g(x)=(8-4x)e^{x}+16$.
\begin{TAB}
\tkzTabInit[lgt=2,espcl=3.5]{$x$/1,$g'(x)$/1,$g(x)$/2}{$-\infty$,$1$,$+\infty$}
\tkzTabLine{,+,z,-,}
\tkzTabVar{-/$16$,+/$g(1)$,-/$-\infty$}
\end{TAB}
أثبت أن المعادلة $g(x)=0$ تقبل حلا وحيدا $\alpha$ حيث $2{,}37<\alpha<2{,}38$ ثم استنتج حسب قيم $x$ إشارة $g(x)$.

\textbf{II)} $f$ الدالة المعرفة على $\R$ بـ: $f(x)=\dfrac{e^{x}+4x}{e^{x}+4}$، و $(C_f)$ تمثيلها البياني في المستوي المنسوب إلى المعلم المتعامد والمتجانس $(O;\vec i,\vec j)$.
\QQ{1} احسب $\lim\limits_{x\to-\infty}f(x)$ ثم بين أن $\lim\limits_{x\to+\infty}f(x)=1$ وفسرها هندسيا.
\QQ{2}
\begin{sq}
\qa بين أنه من أجل كل $x$ من $\R$: $f'(x)=\dfrac{g(x)}{\left(e^{x}+4\right)^2}$.
\qb استنتج اتجاه تغير الدالة $f$، ثم شكل جدول تغيراتها.
\end{sq}
\QQ{3}
\begin{sq}
\qa بين أن المستقيم $(\Delta)$ ذا المعادلة $y=x$ مقارب مائل للمنحنى $(C_f)$ عند $-\infty$.
\qb أدرس الوضع النسبي للمنحنى $(C_f)$ والمستقيم $(\Delta)$.
\end{sq}
\QQ{4}
\begin{sq}
\qa أرسم $(\Delta)$ و $(C_f)$. (نأخذ: $f(\alpha)\approx1{,}4$).
\qb عين بيانيا قيم الوسيط الحقيقي $m$ الموجب تماما حتى تقبل المعادلة $f(x)=\ln(m)$ حلين مختلفين.
\end{sq}
\end{tamrin}

% ===================== بكالوريا 2023 -- الموضوع الأول -- رياضيات =====================
\Needspace{16\baselineskip}
\begin{tamrin}{بكالوريا 2023 -- الموضوع الأول -- رياضيات}
\textbf{I)} $g$ الدالة المعرفة على $\R$ بـ: $g(x)=1+(6x-3)e^{-2x}$.
\QQ{1}
\begin{sq}
\qa أحسب $\lim\limits_{x\to-\infty}g(x)$ و $\lim\limits_{x\to+\infty}g(x)$.
\qb أدرس اتجاه تغير الدالة $g$، ثم شكل جدول تغيراتها.
\end{sq}
\QQ{2}
\begin{sq}
\qa أثبت أن المعادلة $g(x)=0$ تقبل حلا وحيدا $\alpha$ حيث: $0{,}2<\alpha<0{,}3$.
\qb استنتج حسب قيم $x$ إشارة $g(x)$.
\end{sq}

\textbf{II)} $f$ الدالة المعرفة على $\R$ بـ: $f(x)=x+1-3xe^{2x}$، و $(C_f)$ تمثيلها البياني في المستوي المنسوب إلى المعلم المتعامد والمتجانس $(O;\vec i,\vec j)$ (الوحدة $2cm$).
\QQ{1}
\begin{sq}
\qa أحسب $\lim\limits_{x\to-\infty}f(x)$ و $\lim\limits_{x\to+\infty}f(x)$.
\qb بين أن المستقيم $(\Delta)$ ذا المعادلة $y=x+1$ مقارب مائل لـ $(C_f)$ عند $-\infty$.
\qc أدرس وضعية $(C_f)$ بالنسبة إلى $(\Delta)$.
\end{sq}
\QQ{2}
\begin{sq}
\qa بين أنه من أجل كل عدد حقيقي $x$: $f'(x)=g(-x)$.
\qb استنتج أن الدالة $f$ متزايدة تماما على المجال $]{-}\infty;-\alpha]$ ومتناقصة تماما على المجال $[-\alpha;+\infty[$، ثم شكل جدول تغيراتها.
\end{sq}
\QQ{3}
\begin{sq}
\qa أثبت أن $(C_f)$ يقبل مماسا $(T)$ موازيا لـ $(\Delta)$، يطلب تعيين معادلة له.
\qb أرسم $(\Delta)$، $(T)$ و $(C_f)$ على $\left]{-}\infty;\frac12\right]$. (نأخذ: $f(0{,}25)\approx0$، $f(-1{,}3)\approx0$ و $f(-\alpha)\approx1{,}2$).
\qc عين بيانيا قيم الوسيط الحقيقي $m$ التي من أجلها تقبل المعادلة $f(x)=x+m$ حلين بالضبط.
\end{sq}
\end{tamrin}

% ===================== بكالوريا 2022 -- الموضوع الثاني -- رياضيات =====================
\Needspace{16\baselineskip}
\begin{tamrin}{بكالوريا 2022 -- الموضوع الثاني -- رياضيات}
$f$ الدالة المعرفة على $]{-}\infty;1[$ كما يلي: $f(x)=\dfrac{e^{x}-x^2}{x-1}$، و $(C)$ تمثيلها البياني في المستوي المنسوب إلى المعلم المتعامد والمتجانس $(O;\vec i,\vec j)$.
\QQ{1} أحسب $\lim\limits_{x\to-\infty}f(x)$ و $\lim\limits_{x\to1^{-}}f(x)$.
\QQ{2}
\begin{sq}
\qa بين أنه من أجل كل عدد حقيقي $x$: $e^{x}-x>0$.
\qb بين أنه من أجل كل عدد حقيقي $x$ من $]{-}\infty;1[$: $f'(x)=\dfrac{(x-2)\left(e^{x}-x\right)}{(x-1)^2}$.
\qc استنتج اتجاه تغير الدالة $f$، ثم شكل جدول تغيراتها.
\end{sq}
\QQ{3}
\begin{sq}
\qa أحسب $\lim\limits_{x\to-\infty}\big[f(x)+x\big]$، ثم فسر النتيجة بيانيا.
\qb أدرس وضعية $(C)$ بالنسبة إلى المستقيم $(\Delta)$ ذي المعادلة $y=-x-1$.
\end{sq}
\QQ{4} أكتب معادلة للمستقيم $(T)$ مماس المنحنى $(C)$ في النقطة ذات الفاصلة $0$.
\QQ{5}
\begin{sq}
\qa بين أن المعادلة $f(x)=0$ تقبل حلا وحيدا $\alpha$ حيث $-0{,}8<\alpha<-0{,}7$.
\qb أنشئ $(\Delta)$، $(T)$ و $(C)$.
\end{sq}
\QQ{6} ناقش بيانيا، حسب قيم الوسيط الحقيقي $m$، عدد وإشارة حلول المعادلة: $\dfrac{e^{x}-x^2+x-1}{x-1}=mx$.
\QQ{7} $g$ الدالة المعرفة على $]{-}\infty;1[$ بـ: $g(x)=\dfrac{\left|e^{x}-x^2\right|}{x-1}$، و $(C_g)$ تمثيلها البياني في نفس المعلم السابق.
أكتب $g(x)$ دون رمز القيمة المطلقة، ثم أنشئ $(C_g)$.
\end{tamrin}

% ===================== بكالوريا 2021 -- الموضوع الأول -- رياضيات =====================
\Needspace{16\baselineskip}
\begin{tamrin}{بكالوريا 2021 -- الموضوع الأول -- رياضيات}
\textbf{I)} $g$ الدالة العددية المعرفة على $\R$ بـ: $g(x)=\left(x^2-3\right)e^{x}+3$.
\QQ{1} أدرس تغيرات الدالة $g$ ثم شكل جدول تغيراتها.
\QQ{2}
\begin{sq}
\qa بين أن المعادلة $g(x)=0$ تقبل في المجال $[1;+\infty[$ حلا وحيدا $\alpha$ يحقق: $1{,}53<\alpha<1{,}54$.
\qb أحسب $g(0)$ ثم استنتج حسب قيم العدد الحقيقي $x$ إشارة $g(x)$.
\end{sq}

\textbf{II)} $f$ الدالة العددية المعرفة على $\R$ بـ: $f(x)=3x+1+\left(x^2-2x-1\right)e^{x}$، و $(C_f)$ تمثيلها البياني في المستوي المنسوب إلى المعلم المتعامد المتجانس $(O;\vec i,\vec j)$.
\QQ{1} أحسب $\lim\limits_{x\to-\infty}f(x)$ و $\lim\limits_{x\to+\infty}f(x)$.
\QQ{2}
\begin{sq}
\qa بين أنه من أجل كل عدد حقيقي $x$: $f'(x)=g(x)$.
\qb استنتج أن الدالة $f$ متزايدة تماما على كل من $]{-}\infty;0]$ و $[\alpha;+\infty[$ ومتناقصة تماما على $[0;\alpha]$.
\qc شكل جدول تغيرات الدالة $f$.
\end{sq}
\QQ{3}
\begin{sq}
\qa بين أن المستقيم $(\Delta)$ ذا المعادلة $y=3x+1$ مقارب مائل لـ $(C_f)$ عند $-\infty$.
\qb أدرس وضعية $(C_f)$ بالنسبة لـ $(\Delta)$.
\qc بين أن $(C_f)$ يمس حامل محور الفواصل في المبدأ $O$، ويقطعه في نقطة وحيدة أخرى فاصلتها $\beta$ تحقق: $2{,}03<\beta<2{,}04$.
\qd بين أن $(C_f)$ يقبل مماسين $(T)$ و $(T')$ موازيين لـ $(\Delta)$. (لا يطلب كتابة معادلة لـ $(T)$ و $(T')$).
\end{sq}
\QQ{4} أرسم $(\Delta)$، $(T)$، $(T')$ و $(C_f)$ على $\left]{-}\infty;1+\sqrt2\right]$.
(نأخذ: $\alpha\approx1{,}53$، $f(\alpha)\approx-2{,}3$، $f(\sqrt3)\approx-2{,}1$ و $f(-\sqrt3)\approx-3{,}2$).
\QQ{5} $h$ الدالة العددية المعرفة على المجال $]0;+\infty[$ بـ: $h(x)=f(\ln x)$.
\begin{sq}
\qa أحسب $\lim\limits_{x\to0^{+}}h(x)$ و $\lim\limits_{x\to+\infty}h(x)$.
\qb أدرس اتجاه تغير الدالة $h$ ثم شكل جدول تغيراتها.
\end{sq}
\end{tamrin}

\begin{molahada}
في نص المصدر: «بين أن المعادلة $g(x)=0$ تقبل حلا وحيدا $\alpha$» و«بين أن $(C_f)$ يقطع حامل محور الفواصل في نقطة وحيدة فاصلتها $\beta$». لكن $g(0)=0$ (فالمعادلة $g(x)=0$ تقبل حلين: $0$ و $\alpha$)، و $f(0)=0$ مع $f'(0)=0$ (فالمنحنى يمس محور الفواصل في $O$). صححنا السؤالين: الحل $\alpha$ وحيد في $[1;+\infty[$، و $\beta$ هي نقطة التقاطع الوحيدة الأخرى.
\end{molahada}

% ===================== بكالوريا 2020 -- الموضوع الثاني -- رياضيات =====================
\Needspace{16\baselineskip}
\begin{tamrin}{بكالوريا 2020 -- الموضوع الثاني -- رياضيات}
\textbf{I)} الدالتان $g$ و $h$ معرفتان على المجال $]{-}\infty;0]$ كما يلي: $g(x)=-2e^{x}$ و $h(x)=x\left(e^{x}+1\right)$.

حدد إشارة كل من $g(x)$ و $h(x)$ على المجال $]{-}\infty;0]$.

\textbf{II)} الدالة العددية $f$ معرفة على المجال $]{-}\infty;0]$ كما يلي: $f(x)=(x-3)e^{x}+\dfrac12x^2$، و $(C_f)$ تمثيلها البياني في المستوي المنسوب إلى المعلم المتعامد والمتجانس $(O;\vec i,\vec j)$.
\QQ{1}
\begin{sq}
\qa بين أنه من أجل كل عدد حقيقي $x$ من المجال $]{-}\infty;0]$: $f'(x)=h(x)+g(x)$.
\qb استنتج اتجاه تغير الدالة $f$ على المجال $]{-}\infty;0]$.
\end{sq}
\QQ{2} أحسب $f(0)$ و $\lim\limits_{x\to-\infty}f(x)$، ثم شكل جدول تغيرات الدالة $f$.
\QQ{3} بين أن المعادلة $f(x)=0$ تقبل حلا وحيدا $\alpha$ في المجال $]{-}\infty;0]$، ثم تحقق أن: $-1{,}5<\alpha<-1{,}4$.
\QQ{4} $(P)$ هو التمثيل البياني للدالة $x\mapsto\dfrac12x^2$ على المجال $]{-}\infty;0]$.
\begin{sq}
\qa أحسب $\lim\limits_{x\to-\infty}\left[f(x)-\dfrac12x^2\right]$، ثم فسر النتيجة بيانيا.
\qb أدرس الوضع النسبي للمنحنيين $(C_f)$ و $(P)$.
\qc أنشئ $(P)$ ثم المنحنى $(C_f)$ على المجال $]{-}\infty;0]$.
\end{sq}
\QQ{5} ليكن $m$ وسيطا حقيقيا. ناقش بيانيا وحسب قيم $m$ عدد حلول المعادلة $|f(x)|=e^{m}$ في المجال $]{-}\infty;0]$.
\end{tamrin}

% ===================== بكالوريا 2019 -- الموضوع الثاني -- رياضيات =====================
\Needspace{16\baselineskip}
\begin{tamrin}{بكالوريا 2019 -- الموضوع الثاني -- رياضيات}
\textbf{I)} $f_k$ الدالة العددية المعرفة على $\R$ كما يلي: $f_k(x)=(x+1)^2e^{-kx}$، و $k$ وسيط حقيقي. $(C_k)$ المنحنى الممثل للدالة $f_k$ في المستوي المنسوب إلى المعلم المتعامد والمتجانس $(O;\vec i,\vec j)$.
\QQ{1} بين أن كل المنحنيات $(C_k)$ تمر من نقطتين ثابتتين يطلب تعيينهما.
\QQ{2} أحسب نهايتي الدالة $f_k$ عند $-\infty$ وعند $+\infty$. (ناقش حسب قيم الوسيط الحقيقي $k$).
\QQ{3}
\begin{sq}
\qa أحسب $f_k'(x)$، ثم حدد حسب قيم الوسيط الحقيقي $k$ اتجاه تغير الدالة $f_k$.
\qb شكل جدول تغيرات الدالة $f_k$ من أجل $k$ عدد حقيقي موجب تماما.
\end{sq}
\QQ{4} ناقش حسب قيم الوسيط الحقيقي $k$ الأوضاع النسبية للمنحنيين $(C_k)$ و $(C_{k+1})$.

\textbf{II)} $f$ الدالة المعرفة على $\R$ بـ: $f(x)=(x+1)^2e^{-2x}$، و $(C_f)$ تمثيلها البياني في المستوي المنسوب إلى المعلم المتعامد والمتجانس $(O;\vec i,\vec j)$.
\QQ{1} شكل جدول تغيرات الدالة $f$، ثم أرسم المنحنى $(C_f)$ على المجال $\left[-\frac32;+\infty\right[$.
\QQ{2}
\begin{sq}
\qa بين أن المعادلة $f(x)=1$ تقبل في $\R$ حلين أحدهما $\alpha$ بحيث: $-1{,}28<\alpha<-1{,}27$.
\qb عين قيم العدد الحقيقي $m$ التي من أجلها تقبل المعادلة $\left|\dfrac{x+1}{e^{x}}\right|=\left|\dfrac{m+1}{e^{m}}\right|$ حلا وحيدا.
\end{sq}
\end{tamrin}

% ===================== بكالوريا 2018 -- الموضوع الثاني -- رياضيات =====================
\Needspace{16\baselineskip}
\begin{tamrin}{بكالوريا 2018 -- الموضوع الثاني -- رياضيات}
\textbf{I)} $g$ الدالة العددية المعرفة على المجال $]0;+\infty[$ كما يلي: $g(x)=\left(1+x+x^2\right)e^{-\frac1x}-1$.
\QQ{1} بين أنه من أجل كل $x$ من المجال $]0;+\infty[$: $g'(x)=\dfrac{(x+1)\left(2x^2+1\right)}{x^2}e^{-\frac1x}$، واستنتج اتجاه تغير الدالة $g$ على المجال $]0;+\infty[$.
\QQ{2} بين أن المعادلة $g(x)=0$ تقبل حلا وحيدا $\alpha$ حيث $0{,}9<\alpha<1$، واستنتج إشارة $g(x)$ على المجال $]0;+\infty[$.

\textbf{II)} $f$ الدالة العددية المعرفة على المجال $]0;+\infty[$ كما يلي: $f(x)=\dfrac1x+(1+x)e^{-\frac1x}$، و $(C_f)$ تمثيلها البياني في المستوي المنسوب إلى المعلم المتعامد والمتجانس $(O;\vec i,\vec j)$.
\QQ{1}
\begin{sq}
\qa أحسب $\lim\limits_{x\to0^{+}}f(x)$ و $\lim\limits_{x\to+\infty}f(x)$.
\qb بين أنه من أجل كل $x$ من المجال $]0;+\infty[$: $f'(x)=\dfrac{g(x)}{x^2}$، واستنتج اتجاه تغير الدالة $f$، ثم شكل جدول تغيراتها.
\end{sq}
\QQ{2} بين أن: $\lim\limits_{x\to+\infty}\left(xe^{-\frac1x}-x\right)=-1$ (يمكن وضع $t=-\frac1x$) ثم استنتج أن المستقيم $(\Delta)$ ذا المعادلة $y=x$ مقارب للمنحنى $(C_f)$ عند $+\infty$.
\QQ{3} $h$ الدالة العددية المعرفة على المجال $]0;+\infty[$ كما يلي: $h(x)=\dfrac1x-1+e^{-\frac1x}$.
\begin{sq}
\qa أحسب $\lim\limits_{x\to+\infty}h(x)$ وأدرس اتجاه تغير الدالة $h$ واستنتج إشارة $h(x)$ على المجال $]0;+\infty[$.
\qb تحقق أن: $f(x)-x=(1+x)h(x)$، ثم استنتج الوضعية النسبية لـ $(C_f)$ بالنسبة للمستقيم $(\Delta)$.
\end{sq}
\QQ{4} أرسم المستقيم $(\Delta)$ والمنحنى $(C_f)$. (نأخذ $f(\alpha)\approx1{,}73$).
\QQ{5} $(u_n)$ متتالية عددية معرفة على $\N^{*}$ بحدها العام $u_n$ حيث: $u_n=\dfrac{n}{n+1}f\!\left(\dfrac1n\right)-\dfrac{n^2}{n+1}$.
\begin{sq}
\qa أكتب $u_n$ بدلالة $n$ ثم بين أن المتتالية $(u_n)$ هندسية يطلب تعيين أساسها وحدها الأول $u_1$.
\qb أحسب بدلالة $n$ المجموع $S_n$ حيث:
\[S_n=\left(\frac12f(1)-\frac12\right)+\left(\frac23f\!\left(\frac12\right)-\frac13\right)+\left(\frac34f\!\left(\frac13\right)-\frac14\right)+\dots+\left(\frac{n}{n+1}f\!\left(\frac1n\right)-\frac1{n+1}\right)\]
\end{sq}
\end{tamrin}

% ===================== بكالوريا 2017 -- الموضوع الثاني -- رياضيات -- الدورة الاستثنائية =====================
\Needspace{16\baselineskip}
\begin{tamrin}{بكالوريا 2017 -- الموضوع الثاني -- رياضيات -- الدورة الاستثنائية}
المستوي منسوب إلى المعلم المتعامد والمتجانس $(O;\vec i,\vec j)$ حيث $\|\vec i\|=1cm$.

\textbf{I)} نعتبر الدالة العددية $f$ المعرفة على $\R$ كما يلي: $f(x)=(x+1)^2e^{-x}$، و $(C)$ تمثيلها البياني.
\QQ{1} أحسب $\lim\limits_{x\to-\infty}f(x)$ و $\lim\limits_{x\to+\infty}f(x)$.
\QQ{2} أدرس اتجاه تغير الدالة $f$ ثم شكل جدول تغيراتها.
\QQ{3} أثبت أن المنحنى $(C)$ يقبل نقطتي انعطاف يطلب تعيين إحداثييهما، أحسب $f(-2)$، ثم أرسم المنحنى $(C)$.

\textbf{II)} ليكن $m$ وسيطا حقيقيا، نعتبر الدالة $f_m$ المعرفة على $\R$ كما يلي: $f_m(x)=\left(x^2+mx+1\right)e^{-x}$، وليكن $(C_m)$ تمثيلها البياني في المعلم السابق.
\QQ{1} أثبت أن جميع المنحنيات $(C_m)$ تشمل نقطة ثابتة $\omega$ يطلب تعيين إحداثييها.
\QQ{2} أدرس اتجاه تغير الدالة $f_m$ واستنتج قيم $m$ التي من أجلها تقبل $f_m$ قيمتين حديتين يطلب تعيينهما.
\QQ{3} $M_m$ نقطة من $(C_m)$ فاصلتها $x_m$ حيث: $x_m=1-m$.
أثبت أنه عندما $m$ يمسح $\R$ فإن $M_m$ تنتمي إلى منحنى يطلب تعيين معادلة له.
\QQ{4} أدرس حسب قيم الوسيط $m$ حيث $m\neq2$ الوضعية النسبية للمنحنيين $(C)$ و $(C_m)$.
\end{tamrin}

% ===================== بكالوريا 2017 -- الموضوع الأول =====================
\begin{tamrin}{بكالوريا 2017 -- الموضوع الأول -- رياضيات}
نعتبر الدالة العددية $f$ المعرفة على $\R$ كما يلي: $f(x)=(-x^3+2x^2)e^{-x+1}$،
و $(C_f)$ تمثيلها البياني في المستوي المنسوب إلى المعلم المتعامد والمتجانس $(O;\vec i,\vec j)$.
\QQ{1}
\begin{sq}
\qa أحسب $\lim\limits_{x\to+\infty}f(x)$ و $\lim\limits_{x\to-\infty}f(x)$، ثم استنتج وجود مستقيم مقارب للمنحنى $(C_f)$ يطلب تعيين معادلة له.
\qb بيّن أنه من أجل كل عدد حقيقي $x$: $f'(x)=x(x^2-5x+4)e^{-x+1}$، ثم استنتج اتجاه تغيّر الدالة $f$ وشكّل جدول تغيّراتها.
\end{sq}
\QQ{2} أكتب معادلة المماس $(T)$ للمنحنى $(C_f)$ في النقطة ذات الفاصلة $2$.
\QQ{3} $h$ الدالة المعرفة على المجال $[0;+\infty[$ كما يلي: $h(x)=x^2e^{-x+2}-4$.
أدرس اتجاه تغيّر الدالة $h$ ثم استنتج إشارة $h(x)$، ثم حدّد عندئذ وضعية المنحنى $(C_f)$ بالنسبة إلى $(T)$ على المجال $[0;+\infty[$.
\QQ{4} أرسم المماس $(T)$ والمنحنى $(C_f)$ على المجال $[0;+\infty[$.
\QQ{5} نعتبر الوسيط الحقيقي $m$ والمعادلة ذات المجهول الحقيقي $x$ الموجب: \quad $f(x)=m(x-2)\ \ \ldots(E)$.
ناقش بيانيا وحسب قيم $m$ عدد حلول المعادلة $(E)$.
\QQ{6} $g$ الدالة العددية المعرفة على المجال $]0;+\infty[$ بـ: $g(x)=f\left(\dfrac1x\right)$.
اعتمادا على السؤال رقم (1)، شكّل جدول تغيّرات الدالة $g$.
\end{tamrin}

% ===================== بكالوريا 2016 -- الموضوع الثاني =====================
\begin{tamrin}{بكالوريا 2016 -- الموضوع الثاني -- رياضيات}
\textbf{I)} $\varphi$ الدالة العددية المعرفة على $\R$ كما يلي: $\varphi(x)=(x^2-x+1)e^{-x+1}-1$.
\QQ{1}
\begin{sq}
\qa أحسب $\lim\limits_{x\to-\infty}\varphi(x)$ و $\lim\limits_{x\to+\infty}\varphi(x)$.
\qb أدرس اتجاه تغيّر الدالة $\varphi$ ثم شكّل جدول تغيّراتها.
\end{sq}
\QQ{2} بيّن أن المعادلة $\varphi(x)=0$ تقبل في $\R$ حلا $\alpha$ يختلف عن $1$، ثم تحقق أن: $2{,}79<\alpha<2{,}80$.
\QQ{3} استنتج إشارة $\varphi(x)$ على $\R$.

\textbf{II)} $f$ و $g$ الدالتان المعرفتان على $\R$ كما يلي: $f(x)=(2x-1)e^{-x+1}$ و $g(x)=\dfrac{2x-1}{x^2-x+1}$،
و $(C_f)$ و $(C_g)$ تمثيلاهما البيانيان في المستوي المنسوب إلى المعلم المتعامد والمتجانس $(O;\vec i,\vec j)$.
\QQ{1}
\begin{sq}
\qa أحسب $\lim\limits_{x\to-\infty}f(x)$ و $\lim\limits_{x\to+\infty}f(x)$.
\qb أدرس اتجاه تغيّر الدالة $f$ ثم شكّل جدول تغيّراتها.
\end{sq}
\QQ{2} بيّن أن للمنحنيين $(C_f)$ و $(C_g)$ مماسا مشتركا $(T)$ في النقطة ذات الفاصلة $1$، ثم جد معادلة له.
\QQ{3} أرسم المماس $(T)$ والمنحنى $(C_f)$.
\QQ{4}
\begin{sq}
\qa بيّن أنه من أجل كل عدد حقيقي $x$: $f(x)-g(x)=\dfrac{(2x-1)\varphi(x)}{x^2-x+1}$.
\qb أدرس إشارة الفرق $f(x)-g(x)$، ثم استنتج الوضع النسبي لـ $(C_f)$ و $(C_g)$.
\end{sq}
\end{tamrin}

% ===================== بكالوريا 2015 -- الموضوع الثاني =====================
\begin{tamrin}{بكالوريا 2015 -- الموضوع الثاني -- رياضيات}
$f$ الدالة المعرفة بـ: $f(0)=0$، ومن أجل كل عدد حقيقي $x$ من المجال $]{-}\infty;0[$: $f(x)=(x-1)e^{\frac1x}$،
و $(C_f)$ تمثيلها البياني في المستوي المنسوب إلى المعلم المتعامد والمتجانس $(O;\vec i,\vec j)$.
\QQ{1} أدرس استمرارية الدالة $f$ عند $0$ من اليسار.
\QQ{2} أحسب $\lim\limits_{x\to0^-}\dfrac{f(x)}{x}$، ثم فسّر النتيجة هندسيا.
\QQ{3}
\begin{sq}
\qa أحسب $\lim\limits_{x\to-\infty}f(x)$.
\qb أدرس اتجاه تغيّر الدالة $f$، ثم شكّل جدول تغيّراتها.
\end{sq}
\QQ{4}
\begin{sq}
\qa بيّن أن: $\lim\limits_{x\to-\infty}\left[f(x)-x\right]=0$.
\qb استنتج أن المنحنى $(C_f)$ يقبل مستقيما مقاربا مائلا $(\Delta)$ عند $-\infty$، يطلب تعيين معادلة له.
\end{sq}
\QQ{5} $g$ الدالة المعرفة على $]{-}\infty;0[$ بـ: $g(x)=\dfrac{f(x)}{x}$.
\begin{sq}
\qa أحسب $\lim\limits_{x\to-\infty}g(x)$.
\qb أدرس اتجاه تغيّر الدالة $g$، ثم شكّل جدول تغيّراتها.
\end{sq}
\QQ{6}
\begin{sq}
\qa بيّن أنه من أجل كل عدد حقيقي $x$ من المجال $]{-}\infty;0[$: $f(x)>x$.
\qb استنتج وضعية المنحنى $(C_f)$ بالنسبة إلى المستقيم $(\Delta)$.
\qc أنشئ المنحنى $(C_f)$.
\end{sq}
\QQ{7} $m$ عدد حقيقي، و $h_m$ الدالة ذات المتغير الحقيقي $x$ المعرفة على المجال $]{-}\infty;0[$ بـ: $h_m(x)=xe^{\frac1x}-mx$.
\begin{sq}
\qa أحسب $h'_m(x)$ حيث $h'_m$ هي الدالة المشتقة للدالة $h_m$.
\qb باستعمال $(C_f)$، ناقش بيانيا وحسب قيم الوسيط الحقيقي $m$، عدد حلول المعادلة: $h'_m(x)=0$.
\end{sq}
\end{tamrin}

% ===================== بكالوريا 2014 -- الموضوع الأول =====================
\begin{tamrin}{بكالوريا 2014 -- الموضوع الأول -- رياضيات}
\textbf{I)} الدالة $g$ معرفة على $\R$ بـ: $g(x)=(2-x)e^{x}-1$.
\QQ{1} أدرس تغيّرات الدالة $g$.
\QQ{2} بيّن أن للمعادلة $g(x)=0$ في $\R$ حلين $\alpha$ و $\beta$ حيث: $-1{,}2<\alpha<-1{,}1$ و $1{,}8<\beta<1{,}9$.
\QQ{3} استنتج إشارة $g(x)$ على $\R$.

\textbf{II)} الدالة $f$ معرفة على $\R$ بـ: $f(x)=\dfrac{e^{x}-1}{e^{x}-x}$،
و $(C_f)$ المنحنى الممثل للدالة $f$ في المستوي المنسوب إلى المعلم المتعامد والمتجانس $(O;\vec i,\vec j)$.
\QQ{1} أحسب نهاية الدالة $f$ عند $-\infty$ و $+\infty$ وفسّر النتيجتين بيانيا.
\QQ{2} بيّن أنه من أجل كل عدد حقيقي $x$: $f'(x)=\dfrac{g(x)}{(e^{x}-x)^2}$، واستنتج اتجاه تغيّر الدالة $f$ ثم شكّل جدول تغيّراتها.
\QQ{3} بيّن أن: $f(\alpha)=\dfrac{1}{\alpha-1}$، واستنتج حصرا للعددين $f(\alpha)$ و $f(\beta)$.
\QQ{4} أحسب $f(1)$، ثم أرسم المنحنى $(C_f)$.
\end{tamrin}

% ===================== بكالوريا 2013 -- الموضوع الثاني =====================
\begin{tamrin}{بكالوريا 2013 -- الموضوع الثاني -- رياضيات}
\textbf{I)} الدالة $g$ معرفة على $\R$ بـ: $g(x)=1+(x^2-1)e^{-x}$.
\QQ{1}
\begin{sq}
\qa أحسب $\lim\limits_{x\to-\infty}g(x)$ و $\lim\limits_{x\to+\infty}g(x)$.
\qb أدرس اتجاه تغيّر الدالة $g$، ثم شكّل جدول تغيّراتها.
\end{sq}
\QQ{2} بيّن أن المعادلة $g(x)=0$ تقبل حلين في $\R$، ثم تحقق أن أحدهما معدوم والآخر $\alpha$ حيث: $-0{,}8<\alpha<-0{,}7$.
\QQ{3} استنتج إشارة $g(x)$، حسب قيم العدد الحقيقي $x$.

\textbf{II)} الدالة $f$ معرفة على $\R$ بـ: $f(x)=x-(x+1)^2e^{-x}$،
و $(C_f)$ منحنى الدالة $f$ في المستوي المنسوب إلى المعلم المتعامد والمتجانس $(O;\vec i,\vec j)$ (وحدة الطول $2cm$).
\QQ{1}
\begin{sq}
\qa أحسب $\lim\limits_{x\to-\infty}f(x)$ و $\lim\limits_{x\to+\infty}f(x)$.
\qb بيّن أن المستقيم $(\Delta)$ ذا المعادلة $y=x$ مقارب مائل للمنحنى $(C_f)$ عند $+\infty$.
\qc أدرس وضعية المنحنى $(C_f)$ بالنسبة إلى المستقيم $(\Delta)$.
\end{sq}
\QQ{2}
\begin{sq}
\qa بيّن أنه من أجل كل عدد حقيقي $x$: $f'(x)=g(x)$.
\qb شكّل جدول تغيّرات الدالة $f$ على $\R$ (نأخذ: $f(\alpha)\approx-0{,}9$).
\end{sq}
\QQ{3}
\begin{sq}
\qa بيّن أن المنحنى $(C_f)$ يقبل مماسين، معامل توجيه كل منهما يساوي $1$، يطلب تعيين معادلة لكل منهما.
\qb مثّل $(\Delta)$ والمماسين والمنحنى $(C_f)$.
\qc ناقش بيانيا، وحسب قيم الوسيط الحقيقي $m$، عدد حلول المعادلة ذات المجهول $x$: $(x+1)^2+me^{x}=0$.
\end{sq}
\end{tamrin}

% ===================== بكالوريا 2012 -- الموضوع الأول =====================
\begin{tamrin}{بكالوريا 2012 -- الموضوع الأول -- رياضيات}
\textbf{I)} $g$ الدالة العددية المعرفة على $\R$ بـ: $g(x)=2-xe^{x}$.
\QQ{1} أدرس تغيّرات الدالة $g$، ثم شكّل جدول تغيّراتها.
\QQ{2} بيّن أن المعادلة $g(x)=0$ تقبل حلا وحيدا $\alpha$ على $\R$، ثم تحقق أن: $0{,}8<\alpha<0{,}9$.
\QQ{3} عيّن، حسب قيم $x$، إشارة $g(x)$.

\textbf{II)} $f$ هي الدالة المعرفة على $\R$ كما يلي: $f(x)=\dfrac{2x+2}{e^{x}+2}$،
و $(C_f)$ تمثيلها البياني في المستوي المنسوب إلى المعلم المتعامد والمتجانس $(O;\vec i,\vec j)$ (وحدة الطول $2cm$).
\QQ{1} بيّن أن $\lim\limits_{x\to+\infty}f(x)=0$، ثم فسّر النتيجة بيانيا.
\QQ{2}
\begin{sq}
\qa أحسب $\lim\limits_{x\to-\infty}f(x)$.
\qb بيّن أن المستقيم $(\Delta')$ ذا المعادلة $y=x+1$ مستقيم مقارب لـ $(C_f)$.
\end{sq}
\QQ{3} أدرس وضعية $(C_f)$ بالنسبة إلى كل من $(\Delta)$ و $(\Delta')$، حيث $(\Delta)$ هو المستقيم ذو المعادلة $y=x$.
\QQ{4}
\begin{sq}
\qa بيّن أنه من أجل كل عدد حقيقي $x$: $f'(x)=\dfrac{2g(x)}{(e^{x}+2)^2}$، ثم استنتج اتجاه تغيّر الدالة $f$.
\qb بيّن أن $f(\alpha)=\alpha$، ثم شكّل جدول تغيّرات الدالة $f$.
\end{sq}
\QQ{5} أرسم $(\Delta)$ و $(\Delta')$ و $(C_f)$.
\QQ{6} ناقش بيانيا، حسب قيم الوسيط الحقيقي $m$، عدد وإشارة حلول المعادلة: $f(x)=f(m)$.
\end{tamrin}

% ===================== بكالوريا 2011 -- الموضوع الأول =====================
\begin{tamrin}{بكالوريا 2011 -- الموضوع الأول -- رياضيات}
نعتبر الدالة العددية $f$ المعرفة على $\R$ كما يلي: $f(x)=(3x+4)e^{x}$،
و $(C_f)$ تمثيلها البياني في المستوي المنسوب إلى المعلم المتعامد والمتجانس $(O;\vec i,\vec j)$.
\QQ{1}
\begin{sq}
\qa أحسب $f'$ و $f''$، ثم برهن بالتراجع أنه من أجل كل عدد طبيعي $n$ غير معدوم فإن: $f^{(n)}(x)=(3x+3n+4)e^{x}$، حيث $f'$ ، $f''$ ، $\ldots$ ، $f^{(n)}$ المشتقات المتتابعة للدالة $f$.
\qb استنتج حل المعادلة التفاضلية: $y''=(3x+16)e^{x}$.
\end{sq}
\QQ{2}
\begin{sq}
\qa بيّن أن: $\lim\limits_{x\to-\infty}f(x)=0$ وفسّر النتيجة هندسيا.
\qb أدرس اتجاه تغيّر الدالة $f$، ثم شكّل جدول تغيّراتها.
\end{sq}
\QQ{3}
\begin{sq}
\qa أكتب معادلة للمماس $(\Delta)$ للمنحنى $(C_f)$ في النقطة $\omega$ التي فاصلتها $-\dfrac{10}{3}$.
\qb بيّن أن $\omega$ هي نقطة انعطاف للمنحنى $(C_f)$.
\qc أرسم $(\Delta)$ و $(C_f)$ على المجال $]{-}\infty;0]$.
\end{sq}
\end{tamrin}

% ===================== بكالوريا 2010 -- الموضوع الأول =====================
\begin{tamrin}{بكالوريا 2010 -- الموضوع الأول -- رياضيات}
\textbf{I)} $g$ الدالة العددية المعرفة على $\R$ بـ: $g(x)=(3-x)e^{x}-3$.
\QQ{1} أدرس تغيّرات الدالة $g$.
\QQ{2} بيّن أن المعادلة $g(x)=0$ تقبل حلين مختلفين أحدهما معدوم والآخر $\alpha$ حيث $2{,}82<\alpha<2{,}83$.
\QQ{3} استنتج إشارة $g(x)$ حسب قيم $x$.

\textbf{II)} $f$ الدالة المعرفة على $\R$ بـ: $\begin{cases} f(x)=\dfrac{x^3}{e^{x}-1}\ ;\ x\neq0\\[4pt] f(0)=0\end{cases}$
و $(C_f)$ تمثيلها البياني في المستوي المنسوب إلى المعلم المتعامد والمتجانس $(O;\vec i,\vec j)$.
\QQ{1} بيّن أن الدالة $f$ تقبل الاشتقاق عند $x_0=0$. أكتب معادلة لـ $(T)$ مماس $(C_f)$ عند المبدأ $O$.
\QQ{2}
\begin{sq}
\qa بيّن أن $\lim\limits_{x\to+\infty}x^3e^{-x}=0$، ثم جد $\lim\limits_{x\to+\infty}f(x)$، وأحسب $\lim\limits_{x\to-\infty}f(x)$.
\qb بيّن أنه من أجل $x\neq0$ فإن: $f'(x)=\dfrac{x^2g(x)}{(e^{x}-1)^2}$.
\qc تحقق أن: $f(\alpha)=\alpha^2(3-\alpha)$، ثم عيّن حصرا له.
\qd أنشئ جدول تغيّرات $f$.
\end{sq}
\QQ{3}
\begin{sq}
\qa أحسب $f(x)+x^3$ واستنتج الوضعية النسبية لـ $(C_f)$ و $(C)$ منحنى الدالة $x\mapsto-x^3$.
\qb بيّن أن $\lim\limits_{x\to-\infty}\left[f(x)+x^3\right]=0$ وفسّر النتيجة هندسيا.
\end{sq}
\QQ{4} أنشئ في نفس المعلم المماس $(T)$ والمنحنيين $(C_f)$ و $(C)$.
\end{tamrin}

% ===================== بكالوريا 2008 -- الموضوع الثاني =====================
\begin{tamrin}{بكالوريا 2008 -- الموضوع الثاني -- رياضيات}
\textbf{I)} $f$ الدالة العددية المعرفة على $\R$ بـ: $f(x)=x-1+\dfrac{4}{e^{x}+1}$،
وليكن $(C_f)$ تمثيلها البياني في المستوي المنسوب إلى المعلم المتعامد والمتجانس $(O;\vec i,\vec j)$.
\QQ{1} أدرس تغيّرات الدالة $f$.
\QQ{2}
\begin{sq}
\qa بيّن أن $(C_f)$ يقبل نقطة انعطاف $w$، وأكتب معادلة مماس $(C_f)$ عند $w$.
\qb أثبت أن $w$ مركز تناظر للمنحنى $(C_f)$.
\end{sq}
\QQ{3}
\begin{sq}
\qa أحسب $\lim\limits_{x\to+\infty}\left[f(x)-(x-1)\right]$ و $\lim\limits_{x\to-\infty}\left[f(x)-(x+3)\right]$.
\qb استنتج أن $(C_f)$ يقبل مقاربين يطلب إعطاء معادلة لكل منهما.
\end{sq}
\QQ{4}
\begin{sq}
\qa بيّن أن $(C_f)$ يقطع حامل محور الفواصل في نقطة وحيدة فاصلتها $x_0$ حيث $x_0\in]-2{,}77;-2{,}76[$.
\qb أحسب $f(1)$ و $f(-1)$ (تدوّر النتائج إلى $10^{-2}$)، ثم أرسم $(C_f)$ ومستقيميه المقاربين.
\end{sq}

\textbf{II)} $g$ الدالة العددية المعرفة على $\R$ بـ: $g(x)=-x+3-\dfrac{4}{e^{x}+1}$، وليكن $(C_g)$ تمثيلها البياني.
\QQ{1}
\begin{sq}
\qa بيّن أنه من أجل كل عدد حقيقي $x$ فإن: $g(x)=f(-x)$.
\qb استنتج أنه يوجد تحويل نقطي بسيط يحوّل $(C_f)$ إلى $(C_g)$.
\end{sq}
\QQ{2} أنشئ $(C_g)$ في نفس المعلم السابق (دون دراسة $g$).
\end{tamrin}

\end{document}
